\documentclass[12pt,a4paper]{article}
\usepackage[margin=2.4cm]{geometry}
\usepackage{amsmath,amssymb,bm,mathtools}
\usepackage{fontspec}
\setmainfont{TeX Gyre Termes}
\setsansfont{TeX Gyre Heros}
\setmonofont{Latin Modern Mono}[Scale=0.9]
\usepackage{graphicx,float,caption,subcaption,booktabs,longtable,array,multirow,tabularx}
\usepackage{xcolor}
\usepackage{setspace}
\usepackage{lineno}
\usepackage{microtype}
\usepackage{xurl}
\usepackage{siunitx}
\usepackage{enumitem}
\usepackage{fancyhdr}
\usepackage{titlesec}
\titleformat{\paragraph}[runin]{\normalfont\itshape}{}{0pt}{}
\titlespacing*{\paragraph}{0pt}{1.1ex plus .2ex}{0.7em}
\usepackage[authoryear,round,semicolon,sort&compress]{natbib}
\definecolor{lk}{RGB}{0,70,140}
\usepackage[colorlinks=true,linkcolor=lk,citecolor=lk,urlcolor=lk,breaklinks=true,pdfborder={0 0 0}]{hyperref}
\usepackage[nameinlink,capitalise]{cleveref}
\sisetup{detect-all,separate-uncertainty=false}
\setstretch{1.5}
\setlength{\parskip}{3pt}
\captionsetup{font=small,labelfont=bf,justification=centering,singlelinecheck=true}
\titleformat{\section}{\large\bfseries}{\thesection.}{0.5em}{}
\titleformat{\subsection}{\normalsize\bfseries}{\thesubsection}{0.5em}{}
\titleformat{\subsubsection}{\normalsize\itshape}{\thesubsubsection}{0.5em}{}
\pagestyle{fancy}\fancyhf{}\fancyfoot[C]{\small\thepage}\renewcommand{\headrulewidth}{0pt}
\setlength{\headheight}{14pt}
\bibliographystyle{plainnathl}
\newcommand{\doihl}[1]{doi: \href{https://doi.org/#1}{\nolinkurl{#1}}}
\newcommand{\dd}{\mathrm{d}}
\newcommand{\ET}{\mathit{ET}}
\newcommand{\Ent}{\mathit{ET}_{\mathrm{nt}}}
\newcommand{\Ks}{K_{\mathrm{s}}}
\newcommand{\Se}{S_{\mathrm{e}}}
\newcommand{\TAW}{\mathit{TAW}}
\newcommand{\repourl}{https://github.com/halilunaziru73-creator/Coupled-Water-Fluxes-Pore-Scale-Dynamics-Precision-Irrigation-Mulching-Feedback-Loops}
\newcommand{\repo}{\href{\repourl}{repository}}
\newcommand{\datalink}{\href{\repourl}{Data}}
\newcommand{\codelink}{\href{\repourl/tree/main/Manuscripts}{Code}}
\newcommand{\softlink}{\href{\repourl/tree/main/Software/SWAFL_Desktop}{Software}}
\newcommand{\manuallink}{\href{\manualurl}{Manual}}
\newcommand{\downloadlink}{\href{\exeurl}{SWAFL.exe}}
\newcommand{\figdir}{../figures}
\newcommand{\manualurl}{https://github.com/halilunaziru73-creator/Coupled-Water-Fluxes-Pore-Scale-Dynamics-Precision-Irrigation-Mulching-Feedback-Loops/blob/main/Software/SWAFL_Desktop/docs/SWAFL_Manual.pdf}
\newcommand{\exeurl}{https://github.com/halilunaziru73-creator/Coupled-Water-Fluxes-Pore-Scale-Dynamics-Precision-Irrigation-Mulching-Feedback-Loops/releases/latest/download/SWAFL.exe}
\newcommand{\authorblock}{%
\begin{center}
{\large\textbf{Naziru Halilu}}$^{1,2,*}$\\[3pt]
{\small $^{1}$ Agricultural University of Athens, Athens, Greece\\
$^{2}$ Public University of Navarre, Pamplona, Spain\\
ORCID: \href{https://orcid.org/0009-0006-2324-9599}{0009-0006-2324-9599}\\
$^{*}$ Correspondence: \href{mailto:halilunaziru73@gmail.com}{halilunaziru73@gmail.com}}
\end{center}}
\newcommand{\codeavail}{%
\datalink{} (the complete Daisy ensemble outputs), \codelink{} (analysis and figure scripts) and \softlink{} (the SWAFL Desktop source) are in the project repository; the \manuallink{} (93 pages) can be reviewed or downloaded as a PDF. The dataset (all 10\,190 runs, raw and derived) is archived on Zenodo, \doihl{10.5281/zenodo.23043626}, and the ensemble table used here is also stored in the repository as \texttt{Manuscripts/reproducibility/data/daisy\_ensemble.csv.gz}. The folder \texttt{Manuscripts/reproducibility} is a self-contained package: pinned software versions (Python 3.13.16, NumPy 2.5.3, SciPy 1.18.1, pandas 3.0.5, scikit-learn 1.9.1, Matplotlib 3.11.2, Node.js 22.22.0), fixed random seeds in every stochastic step, SHA-256 checksums, and one command (\texttt{run\_all.sh}) that regenerates all results and figures from the ensemble and then compares every regenerated value with the archived result and with the headline numbers printed in the manuscripts. Each table and figure names the script that produced it (Appendix~\ref{app:repro}). Daisy itself was not re-run for this package, no field measurements are used, and three small intermediate arrays and the interface screenshots are archived rather than regenerated; these limits are listed in \texttt{README\_REPRODUCE.md}.}

\usepackage{pdflscape}
\setlength{\emergencystretch}{3em}

\hypersetup{pdftitle={Soil Hydraulic Memory, the Recoverability of Evaporation Savings and the Closure of the Water Balance: A Simulation Study of 10 190 Daisy Runs},pdfauthor={Naziru Halilu}}
\begin{document}
\linenumbers
\begin{titlepage}
\vspace*{-1.2cm}
\begin{center}\begin{minipage}{0.92\textwidth}\centering{\Large\bfseries Soil Hydraulic Memory, the Recoverability of Evaporation Savings and the Closure of the Water Balance: A Simulation Study of 10\,190 Daisy Runs\par}\end{minipage}\end{center}
\vspace{0.4cm}
\authorblock
\vspace{0.4cm}

{\small\setlength{\parindent}{0pt}\setlength{\parskip}{0pt}
\begin{center}\textbf{\normalsize Abstract}\end{center}\vspace{-0.4em}
\noindent Water that is kept from evaporating is routinely counted as water gained, and the outputs of process models that support this count are treated as closed accounts. Both premises are tested on one designed ensemble of 10\,190 simulations of the Daisy soil--plant--atmosphere model (ten soils, 17 experiments). First, the recoverable share of an evaporation saving is made a measured property of the soil: hydraulic memory, the share of released antecedent storage that drains, runs from 0.02 to 0.85 across the soils, follows the pore-size parameter (Spearman 0.95) and saturated conductivity (0.97), transfers to independent evaporation and mulch experiments (0.85, 0.89), and is predicted for soils held out of the fit by a rate-competition law with one parameter (RMSE 0.058). Warming raises and wetter climates lower the recoverable share; nitrogen, a negative control, has no effect. Second, the accounting that these quantities rest on is audited: a blind linear-invariant search with a permutation null recovers the water balance to within 5\,\%, and then shows an unattributed evaporative term that averages 164\,mm\,yr$^{-1}$ (30\,\% of evapotranspiration) and a storage gap of 16\,\% in the two initial-moisture experiments. A compact three-term law predicts evapotranspiration in withheld experiments and soils as well as gradient boosting ($R^2$ 0.90), whereas drainage resists every law tried. The two parts are joined: judged against the reported soil evaporation alone, the drainage share of an evaporation change differs by more than 0.1 in six of ten soils and changes sign in two, so a saving can be judged only against the closed non-transpiration term. The paper gives a computed, soil- and climate-specific discount for evaporation savings, a template that tells a modeller which outputs can be trusted, and testable predictions for lysimeters. All findings concern a one-dimensional matrix-flow model; confirmation needs measured drainage.

\medskip\noindent \textbf{Keywords.} soil hydraulic memory; evaporation savings; drainage partitioning; water-balance closure; equation discovery; process-model auditing\par}
\end{titlepage}

\section*{List of symbols}
\begin{table}[H]\centering\small
\caption{List of symbols, abbreviations and units used in the text.}\label{tab:notation}
\begin{tabularx}{\textwidth}{>{\raggedright\arraybackslash}p{3.3cm}Xl}
\toprule
Symbol & Meaning & Unit\\\midrule
$h$, $h_0$ & suction head; initial suction in E9 & cm\\
$\theta,\ \Se$ & volumetric water content; effective saturation & $-$\\
$\alpha,\ n,\ m$ & van~Genuchten parameters, $m=1-1/n$ & cm$^{-1}$, $-$, $-$\\
$K_s,\ K,\ K_r$ & saturated, unsaturated and relative conductivity & cm\,d$^{-1}$, $-$\\
$\lambda$ & pore-connectivity exponent (0.5) & $-$\\
$P_r,\ I,\ D,\ \mathit{ET}$ & rainfall, irrigation, drainage at 1\,m, evapotranspiration & mm\,yr$^{-1}$\\
$T,\ E_s,\ E_c,\ \Ent$ & transpiration, soil evaporation, canopy evaporation, $\mathit{ET}-T$ & mm\,yr$^{-1}$\\
$\bar R$ & storage release, \cref{eq:release} & mm\,yr$^{-1}$\\
$\phi_x$ & memory fate of flux $x$, \cref{eq:phi} & $-$\\
$\psi_x$ & loss fate of flux $x$, \cref{eq:psi} & $-$\\
$\rho_R$ & recoverable fraction, $1-\psi_D$ & $-$\\
$q^\ast$ & effective withdrawal scale in \cref{eq:rate} & cm\,d$^{-1}$\\
$\TAW$ & total available water to 1\,m & mm\\
E1--E17 & numbered experiments of the Daisy ensemble (E9, E16, E7, E11 and E14 are used here) & $-$\\
LOSO, RMSE & leave-one-soil-out validation; root-mean-square error & $-$\\\bottomrule
\end{tabularx}
\end{table}

%==========================================================
\medskip\noindent\textbf{Additional symbols of the accounting audit.}\par
\begin{table}[H]\centering\small
\caption{Symbols, abbreviations and units of the accounting audit.}\label{tab:notationB}
\begin{tabularx}{\textwidth}{>{\raggedright\arraybackslash}p{3.3cm}Xl}
\toprule
Symbol & Meaning & Unit\\\midrule
$P_r,\ I,\ W$ & rainfall, irrigation, water input $P_r+I$ & mm\,yr$^{-1}$\\
$\mathit{ET},\ T,\ E_s,\ E_c,\ R$ & evapotranspiration and components; unattributed term, \cref{eq:R} & mm\,yr$^{-1}$\\
$D,\ \Delta S,\ Q_r$ & drainage at 1\,m, storage change, runoff & mm\,yr$^{-1}$\\
$\varepsilon$ & closure residual, \cref{eq:closure} & $-$\\
$\mathbf F,\ \sigma_k,\ \mathbf v_k$ & flux matrix, singular values and right singular vectors & --\\
$\mathbf c$ & invariant coefficients, \cref{eq:law} & $-$\\
$n,\ \alpha$ & van~Genuchten shape and scale & $-$, cm$^{-1}$\\
$\TAW$ & total available water & mm\\
$R^2_{\mathrm{cv}}$ & cross-validated coefficient of determination, \cref{eq:r2cv} & $-$\\
$\kappa$ & irrigation compensation, \cref{eq:kappa} & $-$\\
E1--E17 & numbered experiments of the Daisy ensemble & $-$\\
LOEO, LOSO & leave-one-experiment-out and leave-one-soil-out validation & $-$\\
SVD, PCA & singular-value decomposition; principal component analysis & $-$\\
GBM & gradient-boosting machine used as the flexible benchmark & $-$\\
DPF, NLF, IE & drainage fraction $D/W$; nitrogen-leaching fraction; irrigation efficiency & $-$\\\bottomrule
\end{tabularx}
\end{table}

%==========================================================
\section{Introduction}
\label{sec:intro}
Evaporation is the largest single outflow of water from non-saturated terrestrial surfaces, and its response to warming is central to projected changes in soil moisture, drought and irrigation demand \citep{Seneviratne2010,Dai2013,Berg2018,Sherwood2014}. At the global scale most terrestrial evapotranspiration passes through plants \citep{Jasechko2013,Schlesinger2014}, but the partitioning is sensitive to hydrologic connectivity, that is, to how readily water is exchanged between storage and the fluxes that deplete it \citep{Good2015}. A warmer atmosphere raises the evaporative demand that the surface can, in principle, impose on the soil \citep{Roderick2002,Brutsaert1998}, and in managed systems this extra demand is met by irrigation or by drawing down storage. Practical interventions, including residue and straw mulches \citep{BalwinderSingh2011}, deficit and scheduled irrigation \citep{Fereres2007,Geerts2009} and soil-sensor-based scheduling \citep{Datta2023}, are justified by the premise that evaporation can be curtailed and the saved water kept for crops.

That premise is less secure than it looks. \citet{Ward2008} and \citet{Grafton2018} showed at basin scale that apparent savings frequently do not appear as increases in the water that is usable, because the water ``saved'' by a technical improvement was previously a return flow or was being consumed elsewhere. The same accounting problem exists inside a single soil profile. If a mulch or a lower evaporation factor suppresses soil evaporation, the water not evaporated has several possible fates: it may remain in the root zone and reduce the irrigation requirement, it may be taken up by roots, or it may percolate below the root zone and be lost to the crop. Only the first two of these constitute a recoverable saving. The partitioning among them is controlled by the soil's ability to retain water against gravity, which depends on the pore-size distribution, on its connectivity, and on the unsaturated hydraulic conductivity at the moisture states the crop and the irrigation scheduler actually visit.

The relevant physics is old. Capillary flow in porous media is described by the Richards equation \citep{Richards1931}, closed with a water-retention function and a relative-conductivity function. The van~Genuchten--Mualem family \citep{vanGenuchten1980,Mualem1976} reduces these to a pore-size-distribution parameter $n$, a scaling parameter $\alpha$ and a saturated conductivity $K_s$, and its parameters can be estimated from texture by pedotransfer functions \citep{Schaap2001,Saxton2006,Weber2024}. Similitude of capillary flow \citep{Miller1956}, tortuosity and connectivity of the pore space \citep{Ghanbarian2013,Blunt2017}, and the roles of macropores and structure \citep{Beven1982,Beven2013,Jarvis2007,Strudley2008,Fatichi2020} are all recognised as controls on the drainage response. What is missing is a compact, dimensionless diagnostic that links these pore-scale and hydraulic controls to the managerial question ``if evaporation changes by one millimetre, how much of that millimetre is recoverable?''

The term \emph{hydraulic memory} is used here for the dependence of the water-balance partitioning on the antecedent storage state of the profile. Soil moisture memory is well established in land--atmosphere studies \citep{Koster2001,Koster2004,Entekhabi1996}, where it describes the persistence of soil-moisture anomalies. A narrower and operationally measurable definition is adopted here: the fraction of the antecedent storage that is released from the root-zone store which leaves as drainage, as opposed to supplying evaporation or transpiration. Because a profile that is wet at the start of the season releases more water than one that starts dry, this fate fraction measures how ``leaky'' the soil store is in the sense of water-balance partitioning. It is related to, but distinct from, the Budyko-type partitioning of annual precipitation \citep{Budyko1974,Milly1994,Zhang2001}, which concerns the long-term mean, and to the ecohydrological theory of soil-moisture-dependent loss functions \citep{Rodriguez2004}.

A few studies have simulated how irrigation and evaporation interact with soil type, and the Daisy soil--plant--atmosphere model has been calibrated and validated widely \citep{Hansen1991,Hansen2012,Holbak2026}, including for evaporation from structured soils \citep{Leuther2024}, freeze--thaw effects on hydraulic properties \citep{Leuther2021,MirandaVelez2023}, residue-induced changes in hydraulic properties \citep{Leuther2026} and mulching \citep{Vuaille2024}. None of these defines a transferable recoverability metric, and none tests whether the metric is organised by a connectivity proxy that can be known in advance for a soil. This study therefore (i) defines the memory and loss fates; (ii) computes them from a systematic virtual experiment spanning ten soils; (iii) tests whether pore-connectivity proxies predict them out of sample; (iv) tests how climate and management modify the recoverable share; and (v) registers the outcome of every hypothesis tested, including the formulations that identify what the memory fate is not.

Three developments make a transferable metric timely. First, irrigation water is increasingly allocated on the basis of expected on-farm savings, so that the premise that an evaporation reduction is a saving carries economic and regulatory weight \citep{Perry2017,Pereira2012}; an erroneous premise directs investment to measures whose benefit is mostly a reallocation of drainage. Second, sensor-based scheduling and digital agriculture tools are proliferating \citep{Datta2023,Garcia2020}, and they are usually parameterised by soil class; a class-level statement of how much of an evaporation change is recoverable could directly improve their advice. Third, process-based models are increasingly used as virtual laboratories, and the surrogate and emulator literature \citep{Razavi2012,Asher2015} requires interpretable, low-dimensional summaries of what the process model does if it is to be trusted by users. The diagnostics proposed here are exactly such summaries: two dimensionless numbers per soil that can be computed from a small designed experiment.

(1) The memory fate $\phi_D$, the loss fate $\psi_D$ and the recoverable fraction $\rho_R=1-\psi_D$ are defined with explicit estimators and uncertainty. (2) Within a mechanistic simulator and ten soils, these diagnostics are shown to be ordered by conductivity-related hydraulic proxies and to transfer across experiments. (3) A one-parameter rate-competition law is provided with out-of-sample validation, and the conditions that bound it are quantified. (4) The shifts in the recoverable share caused by warming and rainfall are quantified, including a continuous climate surface for one soil. (5) A hypothesis register records the test and outcome of every formulation examined, so that the definitions are tied to evidence.

The analysis addresses four linked problems. Soils differ in how much of the released antecedent storage drains, and it is tested whether this memory fate $\phi_D$ follows pore-connectivity proxies. Transferability is examined next: $\phi_D$, measured from initial moisture, is used to predict $\psi_D$ from evaporation-change experiments. A single competition law with one fitted parameter is then evaluated against purely structural predictors under leave-one-soil-out validation. Finally, the response of the recoverable share $1-\psi_D$ to warming, rainfall, irrigation trigger, dose and nitrogen supply is quantified. Section~\ref{sec:fail} reports the alternative formulations that were tested and what each establishes about the definition of the memory fate.


\paragraph{The accounting question.} Conservation of mass is the most reliable organising principle in hydrology, and it is the foundation of every water-balance model from the bucket to the Earth system model. Yet in practice the budget of a managed soil profile rarely closes exactly in the output of a process model: the fluxes that are reported separately (transpiration, soil evaporation, canopy interception) need not add to the evapotranspiration that enters the balance; change in storage is reported or implied but its size relative to the fluxes is seldom audited; and boundary fluxes at the base of a one-dimensional column are, by construction, definitions rather than measurements \citep{Hillel1998,Refsgaard2004}. At the same time, a strand of scientific machine learning has argued that conservation laws, governing equations and low-dimensional invariants can be \emph{discovered} from data, using sparse regression over libraries of candidate terms \citep{Brunton2016,Champion2019,deSilva2020}, symbolic regression \citep{Schmidt2009,Udrescu2020}, neural networks that learn invariants \citep{Iten2020,Liu2021}, and physics-informed learning \citep{Raissi2019,Karniadakis2021}. In hydrology, deep learning has delivered impressive predictive skill \citep{Kratzert2018,Kratzert2019,Shen2018}, with an ongoing debate about what such models reveal about processes \citep{Nearing2021,Reichstein2019} and about whether interpretable models should be preferred for consequential decisions \citep{Rudin2019}. Enforcing conservation analytically in learned emulators is also an active topic \citep{Beucler2021}.

Three gaps limit confidence in discovery claims in soil-water science. First, discovery methods are rarely validated against a case with a known answer before being applied to an unknown. Second, a discovered relation is typically validated within the data from which it was found, without leave-one-\emph{experiment}-out tests that ask whether the relation survives a change of regime, which is the relevant test of a ``law'' \citep{Roberts2017}. Third, negative results (where sparse laws fail to beat black-box benchmarks, or where a presumed constant is not constant) are rarely documented, so the reader cannot estimate the base rate of success \citep{Kirchner2006,Oreskes1994}.

A large designed ensemble of process-model runs, in which the true balance is known to hold up to the model's own storage term, is taken as a laboratory for method evaluation. A four-part pipeline is built (\cref{fig:pipeline}): (1) a linear-invariant search by singular-value decomposition of the flux matrix with a permutation null, used first as a positive control; (2) a unattributed-term search by auditing partitions that should sum to a total; (3) a manifold analysis of water-fate composition by compositional PCA \citep{Aitchison1986,Egozcue2003}; and (4) interpretable sparse regression over a dimensionless library with LOSO and LOEO validation against a gradient-boosting benchmark \citep{Friedman2001,Tibshirani1996,Hastie2009}. One fixed benchmark, cross-validation grouping and stopping rule are used for all targets, and every law that failed is reported.

The paper asks what a blind method can and cannot recover from a flux matrix. A linear-invariant search must first return the water balance with calibrated uncertainty, which sets the credibility of everything that follows. Accounting terms that the model's own output columns leave open are then sought, and the number of independent directions governing the partition of water among fates is counted. Compact, dimensionally sensible laws are tested against a flexible machine-learning model in \emph{unseen experiments and unseen soils}. The last aim is to separate relations that hold across regimes from those that only appear universal.

\clearpage
\newgeometry{left=0.6cm,right=0.6cm,top=0.5cm,bottom=0.5cm}
\begin{landscape}
\nolinenumbers
\begin{figure}[p]\centering
\includegraphics[width=24.6cm]{\figdir/p2/fig01_pipeline_A4.pdf}
\caption{Pipeline for interpretable discovery. The pipeline recovers the known water balance as a positive control, audits partitions, measures the dimension of the partition manifold and fits sparse laws with LOSO/LOEO validation against gradient boosting. Controls (dashed purple boxes) are shown at the bottom and outputs at the top; purple hexagons are processing blocks, the pink stippled box is the data store, peach boxes are outputs and dotted lines are flows of information. The 149 runs excluded as numerical failures are analysed in \cref{sec:failed}.}
\label{fig:pipeline}
\end{figure}
\end{landscape}
\restoregeometry



\paragraph{Study design and levels of evidence.} The unit of observation is the 20-year mean water balance of one simulated soil under one combination of experimental factors, and statements about soils are made at the level of the ten soils and not of the individual runs. Four levels of evidence are kept apart. Computational reproducibility, the regeneration of every reported number from archived code and data, is established. Predictive validity is assessed by leave-one-soil-out prediction and by transfer between independent experimental designs. Mechanistic validity is bounded by the process model, which represents matrix flow in one dimension. Field transferability is not tested, and the conclusions are therefore stated as predictions that lysimeter experiments can confirm or refute. For the accounting audit the runs of an experiment are correlated because they share a design, so validation withholds whole experiments and whole soils, and a flexible benchmark is reported even where it wins.

\paragraph{What is new.} Evaporation savings are normally counted as gains, and the outputs of a process model as closed accounts. This paper tests both on one ensemble. (1) The recoverable share of a saving becomes a measured soil property: the memory fate $\phi_D$, estimated from initial moisture alone, predicts the loss fate $\psi_D$ of independent evaporation-change experiments across soils and climates, and a one-parameter rate-competition law explains the ordering with out-of-sample validation and stated bounds. (2) The accounting is audited before any partition is used as a target: a blind search must return the water balance first, with a permutation null, and then exposes an unattributed term of 30\,\% of evapotranspiration and a storage gap of 16\,\% in two experiments; laws are validated across experiments and soils and not only across random folds. (3) The two are joined: judged against reported soil evaporation instead of the closed non-transpiration term, the recoverable share changes by more than 0.1 in six of ten soils and reverses sign in two (\cref{sec:auditlink}). (4) Every formulation that was tried is kept in a hypothesis register with its outcome, including the failures, and the protocol is given as seven rules that other process models can adopt.

\section{Methods of the memory analysis}
\label{sec:methods}
\subsection{Overview of the virtual-experiment design}
The analysis uses a pre-existing ensemble of 10\,190 simulations (149 failed and were excluded, leaving 10\,041 valid runs) produced with Daisy \citep{Hansen2012,Holbak2026}, organised into 17 experiments (E0--E16) that vary irrigation dose, trigger pressure, sensor depth, method, soil hydraulic dynamics, mulch, initial moisture, rainfall, temperature, evaporation factor, nitrogen rate, depth and amount (\cref{tab:exp}). Ten soils (S01--S10) span the texture triangle from sand to silty clay loam. All runs were driven by the same 20-year Taastrup (Denmark) weather series; unless stated otherwise every reported flux is a 20-year mean in mm\,yr$^{-1}$, which suppresses inter-annual noise but also means that results describe long-term mean partitioning, not event response. Daisy was selected because it resolves a one-dimensional Richards-equation soil column, coupled to a crop module, a mulch/residue layer, and nitrogen dynamics, which allows all components of the water balance (rainfall, irrigation, evapotranspiration split into transpiration, soil evaporation and canopy evaporation, drainage at 1\,m, and change in storage) to be extracted for every run.

\begin{table}[H]\centering\small
\caption{Experiments of the ensemble used in this paper. ``Valid'' counts exclude runs that failed to complete. Experiments marked $\ast$ are the primary ones for the diagnostics defined here.}
\label{tab:exp}
\begin{tabular}{llrp{6.6cm}}
\toprule
ID & Name & Valid runs & Factors varied (per soil unless noted)\\\midrule
E0 & Baselines & 10 & One baseline per soil\\
E1 & Doses & 189 & Irrigation dose\\
E2 & TriggerPressure & 197 & Trigger suction\\
E3 & SensorDepth & 200 & Sensor depth\\
E4 & Methods & 96 & Surface, overhead, subsurface drip\\
E5 & SHP\_Dynamics & 300 & Time-varying soil hydraulic properties\\
E6 & Mediterranean\_Pollution & 193 & Climate/N scenarios\\
E7$^\ast$ & Mulching & 200 & Mulch vapour-flux factor (VFF) $\times$ mulch storage capacity (CAP), 20 cells\\
E8 & Rain\_Intensity\_Moisture & 1000 & Rain intensity $\times$ initial moisture (S01)\\
E9$^\ast$ & Soil\_Moisture\_Rain & 987 & Initial suction $h_0$ (10 levels) $\times$ rain scale $P$ (10 levels)\\
E10 & Amount\_Timing\_Soil & 954 & Irrigation amount and timing\\
E11 & Temp\_Evap\_Rain & 1000 & Temperature shift $\times$ evaporation factor $\times$ rain (S01)\\
E12 & Sensor\_Trigger\_Rain & 1000 & Depth $\times$ trigger $\times$ rain (S01)\\
E13 & N\_Water\_Soil & 975 & Nitrogen rate $\times$ water\\
E14 & Mulch\_Rain & 1000 & Mulch cover $\times$ capacity $\times$ rain (S01)\\
E15 & Method\_Depth\_Amount & 1000 & Method $\times$ depth $\times$ amount (S01)\\
E16$^\ast$ & Edge\_Cases & 740 & $2^6$ factorial per soil ($P$, $T$, $h$, $D$, $N$, $E_p$) and 13 further edge cases per soil\\\bottomrule
\end{tabular}
\end{table}

Daisy solves the one-dimensional Richards equation in a layered profile with the retention and conductivity functions above, couples it to a crop module (light interception, photosynthesis, root growth and water stress), a surface mulch/residue layer that intercepts rain and suppresses soil evaporation, an energy-balance-based potential evapotranspiration, and nitrogen transformations \citep{Hansen1991,Hansen2012,Holbak2026}. The irrigation scheduler in the ensemble applies a fixed dose whenever the suction at a sensor depth exceeds a trigger. Drainage is the flux at 1\,m depth; no groundwater table is modelled so that net capillary rise is possible but is a boundary-condition artefact in a few soils (\cref{tab:e9}). Because the soil columns are one-dimensional, lateral flow, run-on, and spatial variability are absent, and the 20-year mean precludes interpretation at the event scale. Model structural limits are discussed in \citet{Refsgaard2004} and \citet{Kirchner2006}: the emphasis here is on whether the model produces the \emph{right answers for the right reasons}, which the present study can address only for hydraulic mechanisms internal to the model.

\subsection{Soil hydraulic description}
The retention and conductivity functions are the van~Genuchten--Mualem pair \citep{vanGenuchten1980,Mualem1976}:
\begin{align}
S_e(h) &= \bigl[1+(\alpha h)^{n}\bigr]^{-m}, \qquad m = 1-\tfrac{1}{n},\label{eq:vg}\\
K(S_e) &= K_s\, S_e^{\,\lambda}\Bigl[1-\bigl(1-S_e^{1/m}\bigr)^{m}\Bigr]^{2},\label{eq:mualem}
\end{align}
where $h$ is the suction head (cm), $S_e=(\theta-\theta_r)/(\theta_s-\theta_r)$ the effective saturation, $\alpha$ (cm$^{-1}$) and $n$ ($-$) shape parameters, $K_s$ the saturated conductivity (cm\,d$^{-1}$) and $\lambda$ the pore-connectivity (tortuosity) exponent, conventionally $\lambda=0.5$ \citep{Mualem1976}. In the diagnostics below, the relative conductivity $K_r=K/K_s$ is used with $\lambda=0.5$. The shape parameter $n$ is a pore-size-uniformity index: large $n$ denotes a narrow pore-size distribution (sand, rapid desaturation), small $n$ a broad distribution (fine-textured soils) with a long tail of fine pores that retain water at high suction. For the ten soils, effective $(\alpha,n)$ were taken from the profile-averaged values in the project documentation (\cref{tab:soils}). Saturated conductivities $K_s$ were available from the class-mean values of \citet{Carsel1988} for eight of the ten soils; for the two soils (S01 coarse loam, S02 coarse sand) whose class does not map to a single USDA texture class, $K_s$ is treated as unknown and implied by the fitted law (Section~\ref{sec:rate}). \emph{It is important that $K_s$ here is a literature class value and not the Daisy input conductivity of each soil horizon.} This is a limitation that propagates to the rate-competition results.

\begin{table}[H]\centering\small
\caption{Soil set. $n$ and $\alpha$: effective van~Genuchten parameters; $\TAW$: total available water to 1\,m used in the project documentation; $\Se(100)$: effective saturation at 100\,cm suction from \cref{eq:vg}.}
\label{tab:soils}
\begin{tabular}{lllrrrr}
\toprule
ID & Texture & $\alpha$ (cm$^{-1}$) & $n$ & $\TAW$ (mm) & $\Se(100)$ & $\phi_D$ (E9)\\\midrule
S01 & Coarse loam & 0.0961 & 1.353 & 130 & 0.445 & 0.029\\
S02 & Coarse sand & 0.0536 & 1.600 & 105 & 0.356 & 0.265\\
S03 & Sand & 0.1450 & 2.680 & 60 & 0.011 & 0.846\\
S04 & Loamy sand & 0.1240 & 2.280 & 75 & 0.040 & 0.741\\
S05 & Sandy loam & 0.0750 & 1.890 & 95 & 0.166 & 0.643\\
S06 & Loam & 0.0360 & 1.560 & 150 & 0.468 & 0.289\\
S07 & Silt loam & 0.0200 & 1.410 & 220 & 0.688 & 0.147\\
S08 & Sandy clay loam & 0.0590 & 1.480 & 160 & 0.418 & 0.285\\
S09 & Clay loam & 0.0190 & 1.310 & 175 & 0.753 & 0.087\\
S10 & Silty clay loam & 0.0100 & 1.230 & 195 & 0.882 & 0.022\\\bottomrule
\end{tabular}
\end{table}

\clearpage
\newgeometry{left=0.6cm,right=0.6cm,top=0.5cm,bottom=0.5cm}
\begin{landscape}
\nolinenumbers
\begin{figure}[p]\centering
\includegraphics[width=24.6cm]{\figdir/p1/fig01_framework_A4.pdf}
\caption{Conceptual and analytical framework. Antecedent storage released from the root zone can leave as drainage, soil evaporation, transpiration or other evaporation (memory fates $\phi$, estimated from E9). An imposed change in evaporative loss is offset by drainage, irrigation or transpiration (loss fates $\psi$, estimated from E16 and E7). The recoverable fraction of an evaporation saving is $1-\psi_D$. Boxes indicate storages and solid lines flows; dashed black lines are flows out of the system, dotted lines flows of information; dashed purple boxes are drivers of soil hydraulic properties and the imposed climate changes. Pink indicates water, orange the crop, light brown the soil and rust the pore network; the purple hexagon is the rate-competition process and peach boxes are results. The brown line is the soil surface.}
\label{fig:framework}
\end{figure}
\end{landscape}
\restoregeometry


\begin{figure}[H]\centering
\includegraphics[width=\textwidth]{\figdir/p1/fig02_hydraulics.pdf}
\caption{Hydraulic properties of the ten soils from \cref{eq:vg,eq:mualem}. The ordering by shape parameter $n$ is the ordering on which the main results are organised.}
\label{fig:hydraulics}
\end{figure}

For each run the 20-year mean annual water balance is
\begin{equation}
P_r + I = \mathit{ET} + D + \Delta S + Q_r,
\label{eq:wb}
\end{equation}
where $P_r$ is rainfall, $I$ irrigation, $\mathit{ET}=T+E_s+E_c$ evapotranspiration (transpiration $T$, soil evaporation $E_s$, canopy evaporation $E_c$), $D$ drainage across the 1\,m boundary, $\Delta S$ the change in profile storage and $Q_r$ runoff. Over the 20-year runs $\Delta S/W$ has a median of 0.9\,\% of the input $W=P_r+I$ (0.65\,\% excluding E8 and E9, whose initial moisture is deliberately perturbed; 95th percentile 2.3\,\%). The \emph{storage release} is defined
\begin{equation}
\bar R \;=\; \mathit{ET}+D-P_r-I,
\label{eq:release}
\end{equation}
which is positive when the profile supplies water to the fluxes and negative when it accumulates water; it equals $-(\Delta S+Q_r)$ by \cref{eq:wb}. The non-transpirational evapotranspiration is further defined as $\Ent=\mathit{ET}-T$, which in the model is $E_s+E_c$ plus a small unattributed residual of the model's evapotranspiration bookkeeping.

\subsection{Memory fate $\phi$ (experiment E9)}
Experiment E9 imposes, for each soil, an initial suction $h_0$ (ten values from 30 to 3000\,cm) and a rainfall scale $P$ (ten values from 0.5 to 1.4), a $10\times10$ lattice (100 runs per soil; 87 for S03 because of failed runs). Wetter initial states (small $h_0$) release more storage, so $\bar R$ varies along the $h_0$ axis. Writing $x\in\{E_s,T,D,\text{oth}\}$ for the soil-evaporation, transpiration, drainage and other-evaporation fluxes (with $\text{oth}=\mathit{ET}-T-E_s$), the \emph{memory fate} of a flux is the within-$P$ pooled regression slope of that flux on $\bar R$:
\begin{equation}
\phi_x \;=\; \frac{\sum_{P}\sum_{c\in P}(\bar R_c-\overline{\bar R}_P)(x_c-\bar x_P)}{\sum_{P}\sum_{c\in P}(\bar R_c-\overline{\bar R}_P)^2},\qquad \phi_{E_s}+\phi_T+\phi_D+\phi_{\mathrm{oth}}=1 ,
\label{eq:phi}
\end{equation}
where the sums run over lattice cells $c$ of one soil and bars denote means within each rain level $P$. The unity constraint follows from \cref{eq:release} when the irrigation response to $\bar R$ is negligible within the lattice. Pooling within $P$ removes the first-order effect of the rainfall change on $\bar R$ and isolates the response to antecedent state. $\phi_D$ is called the \emph{hydraulic memory} of drainage: the fraction of an additional millimetre of antecedent storage released that is shed as drainage. Confidence intervals are obtained by bootstrapping the 100 cells (1\,000 resamples). A field-capacity-referenced storage capacity is also computed, $\mathrm{cap}_{FC}=\bar R(h_0{=}100\text{\,cm})-\bar R(h_0{=}1500\text{\,cm})$ at $P=1$.

\subsection{Loss fate $\psi$ and recoverable fraction (E16 and E7)}
E16 is a $2^6$ factorial per soil with factors rainfall scale $P\in\{0.5,1.5\}$, temperature shift $T\in\{2,6\}$\,\textdegree C, trigger suction $h\in\{50,2000\}$\,cm, dose $D\in\{1.2,30\}$\,mm, nitrogen rate $N\in\{25,150\}$\,kg\,ha$^{-1}$ and soil-evaporation factor $E_p\in\{0.3,1.2\}$; 13 further edge cases per soil (extremes of method, mulch, rooting depth and sensor depth) complete the experiment and are not used here. For each of the $2^5=32$ combinations of the other factors, the paired difference between $E_p=1.2$ and $E_p=0.3$ is formed. The loss fate of component $x$ is the ratio of summed paired differences:
\begin{equation}
\psi_x \;=\; \frac{\sum_{k}\bigl[x_k(E_p{=}1.2)-x_k(E_p{=}0.3)\bigr]}{\sum_{k}\bigl[\Ent{}_k(E_p{=}1.2)-\Ent{}_k(E_p{=}0.3)\bigr]}\times s_x,
\label{eq:psi}
\end{equation}
with $s_x=-1$ for drainage and transpiration (which fall when evaporation rises) and $s_x=+1$ for irrigation (which rises). Thus $\psi_D$ is the drainage given up, per millimetre of extra non-transpirational evaporation; $\psi_I$ the extra irrigation and $\psi_T$ the transpiration lost. In a closed balance $\psi_D+\psi_I+\psi_T+\psi_S\approx1$, where $\psi_S$ is the storage term. Because the shares sum to approximately one (\cref{tab:closure}), a value $\psi_D>1$ is possible only when another term, here irrigation, changes sign; it occurs in sand (Section~\ref{sec:results}). \emph{Recoverable fraction}: because drainage below the root zone is water the crop cannot use, whereas reduced irrigation or increased transpiration represents water retained in the system, the recoverable fraction is defined
\begin{equation}
\rho_R \;=\; 1-\psi_D .
\label{eq:rec}
\end{equation}
E7 varies the mulch vapour-flux factor (VFF, 1 = bare soil) and mulch storage capacity (CAP, mm) in a lattice of 20 cells per soil. Mulch reduces $\Ent$, and each flux is regressed on $\Ent$ across the 20 cells; $\psi_D=-\mathrm{d}D/\mathrm{d}\Ent$ and $\psi_I=\mathrm{d}I/\mathrm{d}\Ent$. A bootstrap confidence interval over cells is computed. Because E7 acts through the physical mulch layer rather than a bulk evaporation factor, it is an independent realisation of the evaporation-saving mechanism.

\subsection{Rate-competition law}
\label{sec:ratelaw}
To explain $\phi_D$ mechanistically two competing pathways are considered for water in the transmission zone after a wetting event. Gravity-driven percolation proceeds at a rate of order $K$ evaluated at the high-moisture end of the profile, while evaporation and root uptake withdraw water at an effective rate $q^\ast$. If the two processes act in parallel on the same pool, the fraction lost to percolation is
\begin{equation}
\phi_D \;=\; \frac{K_s}{K_s+q^\ast},
\label{eq:rate}
\end{equation}
where $q^\ast$ (cm\,d$^{-1}$) is a single effective withdrawal scale shared by all soils. This is a competing-rates (parallel-sink) closure, not a derivation from the Richards equation; it is the simplest functional form that (i) is bounded in $[0,1]$, (ii) tends to 0 for impermeable and to 1 for freely draining soils, and (iii) has a single free parameter. $q^\ast$ is estimated by nonlinear least squares on the eight soils with literature $K_s$ and validated by leave-one-soil-out (LOSO) prediction. A variant uses the unsaturated conductivity at a fitted characteristic suction $h^\ast$,
$\phi_D=K(h^\ast)/[K(h^\ast)+q]$ (Appendix~\ref{app:hstar}).

\label{sec:theory}
From \cref{eq:wb,eq:release}, any perturbation $\delta$ of the system that leaves rainfall unchanged satisfies $\delta\bar R=\delta\mathit{ET}+\delta D-\delta I$. Dividing by $\delta\bar R$ and splitting $\mathit{ET}=T+E_s+\text{oth}$ gives $1=\phi_T+\phi_{E_s}+\phi_{\rm oth}+\phi_D-\phi_I$ with $\phi_I=\delta I/\delta\bar R$. Within the E9 lattice irrigation responds weakly to initial moisture and the identity reduces to the unity constraint in \cref{eq:phi}. For an evaporation perturbation $\delta E$, the same identity with $\delta\mathit{ET}=\delta\Ent+\delta T$ and a closed balance ($\delta\bar R\approx0$ in steady state) gives
\begin{equation}
\delta\Ent = -\delta T-\delta D+\delta I\;\;\Longrightarrow\;\;\psi_T+\psi_D+\psi_I\approx 1,\quad \psi_I=\frac{\delta I}{\delta\Ent},\ \psi_D=-\frac{\delta D}{\delta\Ent},\ \psi_T=-\frac{\delta T}{\delta\Ent},
\label{eq:psiclosure}
\end{equation}
which is what \cref{tab:closure} verifies numerically (sums 1.00--1.03; 1.13 for S10, where a storage term is evidently non-negligible). \Cref{eq:psiclosure} makes the logic of recoverability explicit: an evaporation change is paid for by three accounts, and drainage is the only one that is lost to the crop.

For $K_s\ll q^\ast$, $\phi_D\approx K_s/q^\ast$ is linear in conductivity; a tenfold lower $K_s$ gives a tenfold lower memory. For $K_s\gg q^\ast$, $1-\phi_D\approx q^\ast/K_s$ and the retained share declines inversely with conductivity. The logistic form in $\log K_s$ has its steepest response at $K_s=q^\ast=71.5$\,cm\,d$^{-1}$, which lies within the range of the sandy soils (\cref{fig:mech}a), and this is why the sandy soils differ most among themselves in $\phi_D$ while the fine soils are compressed near zero.

With $\phi_D$ dimensionless, and the governing variables $K_s$ (L\,T$^{-1}$), $q^\ast$ (L\,T$^{-1}$), $n$, and $\alpha h^\ast$, the Buckingham $\Pi$ theorem \citep{Buckingham1914} implies $\phi_D=F(K_s/q^\ast,\,n,\,\alpha h^\ast)$. \Cref{eq:rate} corresponds to dropping the dependence on $n$ and $\alpha h^\ast$ beyond that carried by $K_s$; the variant in Appendix~\ref{app:hstar} retains a dependence on $(n,\alpha)$ through $K(h^\ast)$. That $K_s$ alone suffices to explain $97\%$ of the variance in eight soils indicates that, in this ensemble, $n$ and $\alpha$ enter chiefly via their covariance with $K_s$ in the texture classes, as expected from pedotransfer relationships \citep{Schaap2001,Saxton2006}.

Four soil-level predictors of $\phi_D$ were compared: (i) $\TAW$; (ii) $\Se(100)$ from \cref{eq:vg}; (iii) $n$; and (iv) $\log_{10}K_s$. For (i)--(iii) all ten soils are used, for (iv) eight. Each predictor enters a one-variable linear regression clipped to $[0,1]$. Performance is the RMSE of LOSO predictions, compared with the ``mean-only'' benchmark. Association among diagnostics uses Spearman $\rho$ with (a) exact permutation $p$ values and (b) percentile bootstrap intervals (10\,000 resamples). Paired effects of climate factors on $\psi_D$ in E16 use per-soil paired differences between factor levels with a two-sided exact sign test across soils. With $n=10$ soils, power is limited; intervals are wide and are reported alongside point estimates \citep{Efron1993,Good2005}.

Analyses used Python with NumPy \citep{Harris2020}, SciPy \citep{Virtanen2020}, scikit-learn \citep{Pedregosa2011} and Matplotlib \citep{Hunter2007}. Figure~\ref{fig:framework} uses a brown, orange, pink, rust and purple scheme; the data figures use the Okabe--Ito colour-blind-safe palette. The ensemble and the scripts are available in the project repository (\datalink, \codelink), consistent with the FAIR principles \citep{Wilkinson2016} and with calls for reproducible hydrology \citep{Hutton2016}.

%==========================================================
\section{Methods of the accounting audit}
\label{sec:methodsB}
\subsection{Data}
The Daisy \citep{Hansen2012,Holbak2026} virtual-experiment ensemble is used, as described in the methods of the memory analysis: 10\,190 runs in 17 experiments (E0--E16), ten soils (S01--S10, van~Genuchten $n=1.23$--2.68, \cref{tab:soilsP2}), 20-year Taastrup weather; 149 runs failed and are excluded (\cref{sec:failed}), and 10\,041 valid runs are analysed. Each run contributes 20-year mean annual rainfall $P_r$, irrigation $I$, evapotranspiration $\mathit{ET}$, transpiration $T$, soil evaporation $E_s$, canopy evaporation $E_c$, drainage at 1\,m $D$, and (derived) crop yield, nitrogen leaching and irrigation efficiency. Experiments differ in structure: single-soil (coarse loam) lattices (E8, E11, E12, E14, E15), all-soil factorials and lattices (E9, E16) and one-factor series (E1--E4). The distributions of the key variables are in \cref{tab:desc}.

\begin{table}[H]\centering\small
\caption{Descriptive statistics of the 10\,041 valid runs (mm\,yr$^{-1}$; yield in Mg\,ha$^{-1}$ not shown).}
\label{tab:desc}
\begin{tabular}{lrrrrr}
\toprule
Variable & Mean & SD & 25\,\% & Median & 75\,\%\\\midrule
Rain $P_r$ & 417 & 143 & 297 & 464 & 464\\
Irrigation $I$ & 180 & 191 & 0 & 178 & 268\\
Water input $W=P_r+I$ & 597 & 251 & 464 & 654 & 722\\
Evapotranspiration $\mathit{ET}$ & 546 & 196 & 386 & 611 & 671\\
Drainage $D$ (1\,m) & 60 & 106 & 3 & 34 & 78\\
Unattributed $R$ (\cref{eq:R}) & 164 & 96 & 103 & 129 & 223\\\bottomrule
\end{tabular}
\end{table}

\begin{table}[H]\centering\small
\caption{Soil parameters used as predictors (effective van~Genuchten values from the project documentation; $\TAW$ in mm to 1\,m).}
\label{tab:soilsP2}
\begin{tabular}{lllrr}
\toprule
ID & Texture & $\alpha$ (cm$^{-1}$) & $n$ & $\TAW$\\\midrule
S01 & Coarse loam & 0.0961 & 1.353 & 130\\
S02 & Coarse sand & 0.0536 & 1.600 & 105\\
S03 & Sand & 0.1450 & 2.680 & 60\\
S04 & Loamy sand & 0.1240 & 2.280 & 75\\
S05 & Sandy loam & 0.0750 & 1.890 & 95\\
S06 & Loam & 0.0360 & 1.560 & 150\\
S07 & Silt loam & 0.0200 & 1.410 & 220\\
S08 & Sandy clay loam & 0.0590 & 1.480 & 160\\
S09 & Clay loam & 0.0190 & 1.310 & 175\\
S10 & Silty clay loam & 0.0100 & 1.230 & 195\\\bottomrule
\end{tabular}
\end{table}

\subsection{Failed runs and why they are excluded}
\label{sec:failed}
Of the 10\,190 runs, 149 (1.5\,\%) are flagged as failed in the ensemble and are excluded from every analysis. Because exclusion that is not random can bias a discovery study, the pattern of failure was examined before the runs were dropped, and the argument is set out in four parts.

\paragraph{Failures are concentrated in three soils and in specific experiments.} Of the 149 failures, 129 (87\,\%) occur in the coarse sand S02, 13 in the sand S03 and 7 in the silty clay loam S10; the other seven soils have none (\cref{fig:failed}a). In S02, 24.9\,\% of the 519 runs fail. In S03 all 13 failures belong to the initial-moisture experiment E9, and in S10 all 7 belong to the nitrogen--water and dose experiments (E13, E1). The three soils lie at the extremes or in the anomalous part of the van~Genuchten range ($n=1.60$ for S02, $n=2.68$ for S03, the steepest retention curve, and $n=1.23$ for S10, the flattest), which is where the unsaturated conductivity function changes most rapidly with suction.

\paragraph{Within a failing soil, failure follows the size of the water pulse.} In S02, the fraction of failed runs rises with the irrigation depth applied per event in experiment E10: none of the 20 runs at 1 and 2\,mm fails, 10\,\% at 3\,mm, 50\,\% at 8--10\,mm and all of the 30 runs at 15, 20 and 30\,mm (\cref{fig:failed}b). It also depends on the rain scale: 69\,\% of the S02 runs at a scale of 1.5 fail (22 of 32) and 12\,\% at 0.5 (5 of 42), against none of the 90 runs at intermediate scales 0.6--1.4 (\cref{fig:failed}c). Runs that share the same weather, the same soil and the same code, and differ only in how much water arrives at once, therefore fail or succeed in a systematic order. This is the behaviour of a state-dependent numerical difficulty and not of a corrupted input, which would hit runs irrespective of their forcing.

\paragraph{The outputs of failed runs are physically impossible.} In valid runs the annual drainage at 1\,m is never below $-17.5$\,mm (median in valid S02 runs, 134\,mm). In failed runs it is negative in 65\,\% of cases and falls to $-882$\,mm (\cref{fig:failed}d), and the ratio of evapotranspiration to water input ($P+I$) ranges from 0.54 to 4.34, with a median of 1.09, so that in some runs more water leaves the profile than entered it. A net upward flux of hundreds of millimetres through the lower boundary of a free-draining column, or evapotranspiration several times the water input, cannot be produced by the physics of the model; it is the signature of a solver that has not converged, in which a mass error is absorbed in the boundary or storage terms.

\paragraph{The alternatives are less consistent with the pattern.} A fault in the weather or in the run definition would not depend on soil and on dose in this ordered way; a failure from over-wetting in the field, such as ponding or runoff, is a physical process that the model is built to represent and would not give negative drainage; and a random job failure would spread evenly over soils. Numerical non-convergence of the Richards-equation solver in a stiff regime (a very permeable, quickly draining column that receives a large input, or a column with an extreme retention curve) accounts for the concentration in soils, the order with water input and the impossible outputs together.

\paragraph{Treatment of the excluded runs.} The failed runs are interpreted as numerical failures of the soil-water solver under stiff conditions, not as information about soil hydrology, and they are excluded for that reason; the interpretation follows from the pattern of failure described above. Two consequences follow. The exclusion removes a quarter of the runs of S02, preferentially in the wettest and driest forcing and at the largest doses, so results for this soil, and any law fitted to it, apply to the part of the factor space in which the solver succeeded. The other nine soils are essentially unaffected (at most 2.5\,\% of runs in S03), and the conclusions that rest on across-soil comparisons were tested for sensitivity to the exclusion by re-analysis. With S02 removed (9\,651 runs) the three-term ET law is re-estimated with coefficients $0.652$ for $\ln W$, $-0.297$ for $n$ and $0.029$ for $\ln(1+I)$, against $0.646$, $-0.299$ and $0.029$ for the full ensemble, with LOEO $R^2=0.897$ and LOSO $R^2=0.904$ (full ensemble 0.898 and 0.904), and the unattributed term $R$ is 30\,\% of ET (163\,mm\,yr$^{-1}$; full ensemble 30\,\% and 164\,mm\,yr$^{-1}$). With S02 and S03 both removed (9\,145 runs) the scores are 0.893 and 0.898 and $R$ is 30\,\% (165\,mm\,yr$^{-1}$). The exclusion therefore does not drive the across-soil results.

\begin{figure}[H]\centering
\includegraphics[width=\textwidth]{\figdir/p2/figS_failed.pdf}
\caption{Pattern of the 149 failed runs. (a) Failed runs per soil, as a percentage of the 519 runs per soil (counts above bars). (b) Failed fraction in S02 against irrigation depth per event in E10. (c) Failed fraction in S02 against rain scale. (d) Annual drainage at 1\,m for valid and failed runs, normalised as densities; negative drainage of several hundred millimetres occurs only in failed runs.}
\label{fig:failed}
\end{figure}

\subsection{Governing balance and the flux matrix}
For each run the mean annual water balance is
\begin{equation}
P_r+I \;=\; \mathit{ET}+D+\Delta S+Q_r,\qquad \mathit{ET}=T+E_s+E_c+R,
\label{eq:wb2}
\end{equation}
where $\Delta S$ is storage change, $Q_r$ runoff, and $R$ is defined as the remainder of ET after the three reported components,
\begin{equation}
R \;=\; \mathit{ET}-T-E_s-E_c .
\label{eq:R}
\end{equation}
Equation~\eqref{eq:R} defines $R$ without ascribing a process. Let $\mathbf F\in\mathbb R^{N\times q}$ be the matrix whose rows are runs and whose columns are the fluxes (rain, irrigation, ET, $T$, $E_s$, $E_c$, $D$). If a linear invariant $\sum_j c_jF_{ij}=0$ holds for all runs, then $\mathbf c$ lies in the null space of $\mathbf F$. Define $W=P_r+I$ and the closure residual
\begin{equation}
\varepsilon \;=\; \frac{W-\mathit{ET}-D}{W}\;\approx\;\frac{\Delta S+Q_r}{W}.
\label{eq:closure}
\end{equation}

\subsection{Linear-invariant search by SVD with a permutation null}
The singular-value decomposition $\mathbf F=\mathbf U\boldsymbol\Sigma\mathbf V^{\top}$ with singular values $\sigma_1\ge\dots\ge\sigma_q$ identifies near-null directions as the right singular vectors $\mathbf v_k$ with small $\sigma_k$. Each column is first divided by its root-mean-square value so that no flux dominates through its units; $\mathbf v_k$ is mapped back to mm units and normalised so that the ET coefficient equals $-1$. Because singular values depend on the scale of the columns, significance is judged against a \emph{column-permutation null}: each column of $\mathbf F$ is permuted independently across runs (destroying cross-column dependence but preserving each marginal), and the smallest singular value under the null is recorded; an invariant is accepted only if $\sigma_{\min}$ lies far below the null value. Uncertainty of the coefficients is estimated by a nonparametric bootstrap over runs (resampling 10\,041 rows; 95\,\% percentile intervals). The test is repeated (a) on the four-flux subsystem $(P_r,I,\mathit{ET},D)$, which is the set that appears in the balance; (b) on the seven-column matrix; and (c) excluding E8 and E9, whose initial states are deliberately away from equilibrium. It is emphasised that this is a \emph{positive control}: the expected answer is $\mathbf c=(1,1,-1,-1)$ up to scale and the pipeline is judged by how closely it recovers it, and by whether it fails in a diagnosable way.

Any decomposition that should sum to a total is audited by computing the residual. The audit covers (a) \cref{eq:R}, reporting $R$, its share $R/\mathit{ET}$ by experiment, and its dependence on drivers by regression of $\ln R$ on candidate drivers (irrigation, soil, mulch capacity, temperature, evaporation factor); and (b) the storage closure $\varepsilon$ of \cref{eq:closure} by experiment. The audit is descriptive: it identifies where the budget does not close, and tests candidate attributions by whether $R$ responds to the drivers that would be expected under each.

The partition of water input into fates (transpiration, soil evaporation, canopy evaporation, unattributed $R$, drainage) is compositional: components are positive and sum to a constant, so Euclidean statistics are inappropriate \citep{Aitchison1986,Egozcue2003}. Writing the composition as $\mathbf x=(x_1,\dots,x_5)$ with $x_i=\text{flux}_i/\sum_j\text{flux}_j$, the centred log-ratio (clr) transform is applied
\begin{equation}
\mathrm{clr}(\mathbf x)_i=\ln x_i-\frac1{5}\sum_{j=1}^5\ln x_j ,
\label{eq:clr}
\end{equation}
and PCA \citep{Wold1987} to the transformed data. Drainage is zero or negative in 14\,\% of runs, so zeros are replaced by a floor $\delta$ (in mm\,yr$^{-1}$) before the closure to unit sum; because the results depend strongly on $\delta$, results are reported for $\delta\in\{10^{-3},0.1,1,5,10\}$. Because log-transformations are inherently sensitive to near-zero values, the sensitivity to the floor is mapped across four orders of magnitude and not fixed at a single value, so that the stability of the drainage axis can be judged independently of the floor selected. The number of components needed to explain the variance is a data-driven estimate of the dimension of the partition manifold.

For a target $y$ (log of a flux), a library $\boldsymbol\Theta$ is built of dimensionless or log-transformed candidates (\cref{tab:lib}). The model is
\begin{equation}
\ln y \;=\; \beta_0+\sum_{k\in\mathcal S}\beta_k\,\theta_k+\epsilon,\qquad \mathcal S\subseteq\{1,\dots,p\},
\label{eq:sparse}
\end{equation}
with the active set $\mathcal S$ chosen by forward selection: at each iteration the candidate that maximises the leave-one-experiment-out cross-validated $R^2$ is added, and the search stops when the improvement is below 0.003. The cross-validated $R^2$ is
\begin{equation}
R^2_{\mathrm{cv}}=1-\frac{\sum_i(\ln y_i-\widehat{\ln y}_i^{(-g(i))})^2}{\sum_i(\ln y_i-\overline{\ln y})^2},
\label{eq:r2cv}
\end{equation}
where $\widehat{\ln y}_i^{(-g(i))}$ is the prediction for run $i$ from a model fitted with all runs of its group $g(i)$ held out. Groups are experiments (LOEO, 17 folds) or soils (LOSO, 10 folds). LOEO is the demanding test: it asks whether a law fitted on 16 experimental designs predicts the 17th, which differs in which factors vary and in what range. The root-mean-square error in mm and the mean absolute percentage error after back-transformation are also reported. The benchmark is a histogram-based gradient-boosting regressor \citep{Friedman2001,Pedregosa2011} with 250 iterations on the whole library, validated under the same groupings. Caveat: because forward selection uses the LOEO score, the reported LOEO $R^2$ of the selected model is mildly optimistic (selection bias); LOSO values, which were not used for selection, provide a partly independent check, and coefficient stability across folds is reported.

\begin{table}[H]\centering\small
\caption{Candidate library $\boldsymbol\Theta$ (log-transformed or dimensionless).}
\label{tab:lib}
\begin{tabular}{lll}
\toprule
Symbol & Definition & Role\\\midrule
$\ln P_r$ & log rainfall & climate input\\
$\ln(1+I)$ & log irrigation & management input\\
$\ln W$ & log water input $P_r+I$ & total supply\\
$\ln E_p$ & log soil-evaporation factor & evaporative demand\\
$\Delta T$ & temperature shift (\textdegree C) & warming\\
$n$ & van~Genuchten shape parameter & pore-size distribution\\
$\ln\TAW$ & log total available water & storage capacity\\
$\ln\alpha$ & log van~Genuchten scaling & air-entry scale\\
mulch & VFF$\times$CAP (mulch lattice) & residue layer\\
$I/W$ & irrigation fraction & managed share\\
$\ln(W/\TAW)$ & dimensionless supply-to-storage & wetness index\\\bottomrule
\end{tabular}
\end{table}

\label{sec:svdtheory}
Suppose the true model is $\mathbf F\mathbf c=\boldsymbol\eta$ with $\boldsymbol\eta$ a small run-wise residual (here $-\Delta S-Q_r$ for $\mathbf c=(1,1,-1,-1)$). The SVD gives the best rank-$(q-1)$ approximation of $\mathbf F$ (Eckart--Young), and the right singular vector $\mathbf v_q$ minimises $\|\mathbf F\mathbf v\|_2$ over unit vectors; the corresponding singular value is $\sigma_q=\|\mathbf F\mathbf v_q\|_2$. Because the columns are all measured with error (here, the model's own discretisation and storage noise), $\mathbf v_q$ is a total-least-squares solution, in which errors in all variables are treated symmetrically; this is preferable to regressing one flux on the others, whose coefficients would be attenuated and depend on the choice of dependent variable. When $\eta$ is uncorrelated with the columns, $\sigma_q^2\approx\|\boldsymbol\eta\|^2$, so $\sigma_q/\sqrt N$ is the root-mean-square residual per run in the units of the scaled columns. Before the decomposition each column is divided by its root-mean-square value (so that no flux dominates by its units), and the singular vector is mapped back to mm units afterwards; the singular values are therefore dimensionless, and for the four-flux subsystem $\sigma_q/\sqrt N=5.9/\sqrt{10041}=0.059$, i.e.\ a residual of about 6\,\% of a typical column scale. A caveat is that the matrix is not mean-centred: centring would remove the constant direction, but a conservation law $\sum c_jF_j=0$ has no intercept, so the uncentred matrix is the correct object.

Permuting each column independently preserves marginal distributions and destroys all linear relations between columns, so the smallest singular value under permutation, $\sigma_q^{\mathrm{perm}}$, measures how small $\sigma_q$ would be if the columns were unrelated but had the same scales. A ratio $\sigma_q/\sigma_q^{\mathrm{perm}}=5.9/33.3=0.18$ for the four-flux matrix shows that the null direction reflects genuine dependence. Because fluxes are positively skewed and share a common experiment structure (for example rainfall is identical across many runs), the permutation null is conservative in one respect (it breaks the experimental clustering) and liberal in another (it does not respect the clustering when permuting). Runs are therefore also bootstrapped and intervals reported.

\begin{enumerate}[nosep,label=\arabic*.]
\item Initialise $\mathcal S=\emptyset$, best score $s^\ast=-\infty$.
\item For each candidate $k\notin\mathcal S$, fit \cref{eq:sparse} on the training groups and compute $R^2_{\mathrm{cv}}$ (\cref{eq:r2cv}) for $\mathcal S\cup\{k\}$ with 17 LOEO folds.
\item Let $k^\ast$ be the maximiser; if $R^2_{\mathrm{cv}}(\mathcal S\cup\{k^\ast\})>s^\ast+0.003$, add $k^\ast$ and update $s^\ast$; otherwise stop.
\item Refit on all data to report coefficients; evaluate LOSO using the same terms.
\end{enumerate}
\Cref{tab:path} shows the path for the ET law: $\ln W$ alone gives $R^2_{\mathrm{cv}}=0.801$, adding $n$ raises it to 0.895 and adding $\ln(1+I)$ to 0.898. The third term improves the score by only 0.003, at the stopping threshold, so a two-term law $\ln\mathit{ET}=f(\ln W,n)$ is almost as good and is the safer summary.

The compensation coefficient $\kappa$ measures the share of a rainfall change that is offset by irrigation. For rainfall change $\Delta P_r$ with other factors fixed,
\begin{equation}
\kappa \;=\; -\frac{\Delta I}{\Delta P_r}.
\label{eq:kappa}
\end{equation}
In single-soil lattices $\mathrm d W/\mathrm dP_r=1-\kappa$ is estimated by ordinary least squares; in E16, paired differences are formed of rain scale 0.5 and 1.5 across the 32 combinations of other factors within each soil, and bootstrap the pooled estimate (500 resamples).

All analyses are scripted in Python \citep{Harris2020,Virtanen2020,Pedregosa2011,Hunter2007}; scripts and intermediate results are in the project repository (\datalink, \codelink) and mapped in Appendix~\ref{app:reproB}. Model evaluation follows the recommendations of \citet{Moriasi2007} and \citet{Bennett2013}; $R^2$ is reported on the log scale unless stated. No test data were used for model selection other than through the declared cross-validation.

%==========================================================
\section{Results of the memory analysis}
\label{sec:results}
\subsection{Antecedent storage release drains differently in different soils}
Across the $10\times10$ E9 lattice the storage release $\bar R$ varies strongly with initial suction (\cref{fig:memory}). In sand (S03, $n=2.68$) the release saturates near $-25$\,mm\,yr$^{-1}$ once $h_0$ exceeds about 100\,cm, whereas in the fine-textured soils (S07, S01, S09, S10) $\bar R$ continues to change across the full range of $h_0$ and the curves show a pronounced knee near 300--800\,cm. The rank order of the panels in \cref{fig:memory} (by decreasing $n$) therefore corresponds to a visible change in the shape of the memory curve from a sharp, early-saturating response to a gradual, extended one.

\begin{figure}[H]\centering
\includegraphics[width=\textwidth]{\figdir/p1/fig03_memory_grid.pdf}
\caption{Storage release $\bar R$ (mm\,yr$^{-1}$) versus initial suction $h_0$ in the E9 lattice, one panel per soil (a--j, ordered by decreasing van~Genuchten $n$). Lines are rain scales $P=0.5$--1.4 (colour bar). Wetter starts (small $h_0$) release more storage; the response saturates early in coarse soils and extends over the full $h_0$ range in fine soils.}
\label{fig:memory}
\end{figure}

The memory fates (\cref{eq:phi}; \cref{fig:fates}a, \cref{tab:fates}) differ by more than an order of magnitude among soils. Drainage takes 85\,\% (95\,\% CI 0.80--0.89) of released antecedent storage in sand, 74\,\% in loamy sand and 64\,\% in sandy loam, but only 2--3\,\% in silty clay loam and coarse loam and 9\,\% in clay loam. In the finer soils transpiration takes most of the released storage ($\phi_T=0.75$--0.84 in S01 and S10); in coarse soils evaporation and transpiration together are small ($\phi_{E_s}+\phi_T=0.15$--0.33 in S03--S05).

\begin{figure}[H]\centering
\includegraphics[width=\textwidth]{\figdir/p1/fig04_fates.pdf}
\caption{(a) Memory fates $\phi$ of released antecedent storage (E9; error bars: bootstrap 95\,\% CI for drainage). (b) Loss fates $\psi_D$ (drainage) and $\psi_I$ (irrigation) of an evaporation-factor increase (E16). (c) Drainage share of the evaporation saving from mulch (E7). Soils are ordered by decreasing $\phi_D$.}
\label{fig:fates}
\end{figure}

\begin{table}[H]\centering\small
\caption{Memory fates (E9) and loss fates (E16, E7) by soil. $\phi_D,\phi_T,\phi_{E_s}$: drainage, transpiration and soil-evaporation shares of released antecedent storage; $\psi_D$, $\psi_I$: drainage and irrigation shares of an evaporation change; $\rho_R^{16}=1-\psi_D^{E16}$.}
\label{tab:fates}
\begin{tabular}{lrrrrrrrrr}
\toprule
Soil & $\phi_D$ & $\phi_T$ & $\phi_{E_s}$ & $\mathrm{cap}_{FC}$ (mm) & $\psi_D^{E16}$ & $\psi_I^{E16}$ & $\rho_R^{16}$ & $\psi_D^{E7}$ & $\psi_I^{E7}$\\\midrule
S01 & 0.029 & 0.750 & 0.145 & 82.8 & 0.295 & 0.619 & 0.705 & 0.298 & 0.583\\
S02 & 0.265 & 0.419 & 0.193 & 104.0 & 0.193 & 0.658 & 0.807 & 0.607 & 0.339\\
S03 & 0.846 & 0.080 & 0.074 & 4.2 & 1.210 & $-0.384$ & $-0.210$ & 0.706 & 0.000\\
S04 & 0.741 & 0.153 & 0.100 & 13.2 & 0.799 & 0.219 & 0.201 & 0.687 & $-0.023$\\
S05 & 0.643 & 0.204 & 0.128 & 41.3 & 0.582 & 0.415 & 0.418 & 0.663 & $-0.027$\\
S06 & 0.289 & 0.582 & 0.040 & 131.8 & 0.469 & 0.433 & 0.531 & 0.521 & 0.379\\
S07 & 0.147 & 0.716 & 0.042 & 208.1 & 0.362 & 0.528 & 0.638 & 0.396 & 0.605\\
S08 & 0.285 & 0.555 & 0.078 & 57.0 & 0.533 & 0.386 & 0.467 & 0.549 & 0.332\\
S09 & 0.087 & 0.774 & 0.027 & 167.5 & 0.312 & 0.645 & 0.688 & 0.426 & 0.584\\
S10 & 0.022 & 0.842 & 0.025 & 173.0 & 0.310 & 0.838 & 0.690 & 0.414 & 0.629\\\bottomrule
\end{tabular}
\end{table}

\subsection{Loss fates and the share of an evaporation change that is drainage}
In E16 the drainage share of the extra evaporation $\psi_D$ ranges from 0.19 (coarse sand) to 1.21 (sand; \cref{tab:fates}). In sand a millimetre of additional non-transpirational evaporation removes more than a millimetre of drainage ($\psi_D=1.21$) while irrigation changes in the opposite direction ($\psi_I=-0.38$) and transpiration falls by 0.18; the three terms still sum to 1.00 (\cref{tab:closure}). The cause was not isolated of the negative irrigation response in sand; it may reflect the interaction of the suction trigger with the very small $\TAW$ (60\,mm), but this is a conjecture, and it is the reason $\psi_D>1$ is not interpreted as a literal fraction. In contrast, in silty clay loam (S10) only 31\,\% of the evaporation change is drainage and 84\,\% is offset by irrigation. In the mulch lattice E7, $\psi_D$ ranges from 0.30 (S01) to 0.71 (S03) and falls monotonically from sandy to fine-textured soils (\cref{fig:fates}c), with the exception of the coarse sand S02, where mulch effects on drainage are larger than the E16 contrast suggests (0.61 versus 0.19). The recoverable fraction $\rho_R=1-\psi_D$ is therefore small in sandy soils (0.20--0.42 for S04--S05; $\leq0$ for S03) 0.64--0.71 in the coarse loam and the fine soils (S01, S07, S09, S10), 0.47--0.53 in the loam and sandy clay loam, and 0.81 in the coarse sand (\cref{tab:fates}).

\subsection{Hydraulic proxies of pore connectivity predict the memory fate}
$\phi_D$ is strongly ordered by $n$ ($\rho=0.95$, permutation $p=1.5\times10^{-4}$, bootstrap CI 0.69--1.00), by $\Se(100)$ ($\rho=-0.85$), by $\TAW$ ($\rho=-0.77$) and by $\log_{10}K_s$ over the eight soils with literature $K_s$ ($\rho=0.97$; \cref{fig:predictors}). The ordering by $n$ is consistent with the idea that a narrow pore-size distribution implies a strongly connected network of large pores that drains quickly once the profile is wetted, while the broad distribution of fine soils provides a continuum of small pores that retain water at the tension states visited by the crop. Because $n$ is also correlated with $K_s$ in texture-based datasets \citep{Carsel1988,Schaap2001,Weber2024}, the two rankings cannot be separated here (Section~\ref{sec:disc}).

\begin{figure}[H]\centering
\includegraphics[width=0.92\textwidth]{\figdir/p1/fig05_predictors.pdf}
\caption{Relationship between the memory fate $\phi_D$ (E9) and four soil-level predictors: (a) van~Genuchten $n$, (b) $\log_{10}K_s$ (eight soils), (c) $\TAW$ to 1\,m, (d) effective saturation at 100\,cm. Spearman $\rho$ given in each panel.}
\label{fig:predictors}
\end{figure}

\subsection{Rate-competition law and out-of-sample skill}
\label{sec:rate}
Fitting \cref{eq:rate} to the eight soils with literature $K_s$ gives $q^\ast=71.5$\,cm\,d$^{-1}$ with $R^2=0.971$ (\cref{fig:mech}a). With eight points the parameter is only moderately well determined: a bootstrap over soils (5\,000 resamples, of which 4\,999 contained at least three distinct soils and were used) gives a 95\,\% interval of 58--105\,cm\,d$^{-1}$, and omitting each soil in turn gives values between 66 and 80\,cm\,d$^{-1}$, so the shape of the law is more secure than the exact value of $q^\ast$. Leave-one-soil-out RMSE is 0.058 on the eight soils with literature $K_s$. On the same eight soils the linear models give 0.100 for $n$, 0.064 for $\Se(100)$ and 0.088 for $\TAW$, and the mean-only benchmark gives 0.338; on all ten soils the corresponding values are 0.097, 0.136, 0.183 and 0.318 (\cref{fig:mech}c; \cref{tab:loso}). On like-for-like soils the $K_s$ law therefore reduces the out-of-sample error by 83\,\% relative to the mean and by 42\,\% relative to $n$, but by only 9\,\% relative to $\Se(100)$, so the advantage of the law over the best simple alternative is small and is not distinguishable with eight points; its value lies in the mechanism and in having a single parameter. Two cautions apply. First, the comparison on ten soils involves a different soil set from the $K_s$ model; an earlier logistic fit in $n$ on all ten soils gave an LOO RMSE of 0.111. Second, $q^\ast=71.5$\,cm\,d$^{-1}$ is orders of magnitude larger than the mean evapotranspiration rate ($\approx0.2$\,cm\,d$^{-1}$); it is an effective scale for the competition between drainage and withdrawal at near-saturated states and should not be read as a physical flux. The law implies $K_s\approx2.1$\,cm\,d$^{-1}$ for S01 and $\approx25.8$\,cm\,d$^{-1}$ for S02, which are plausible for a coarse loam and a coarse sand respectively but remain untested predictions (the markers in \cref{fig:mech}a).

\begin{figure}[H]\centering
\includegraphics[width=\textwidth]{\figdir/p1/fig06_mechanism.pdf}
\caption{Rate-competition law. (a) $\phi_D$ against $K_s$ with the fitted curve $K_s/(K_s+q^\ast)$; crosses mark implied $K_s$ for S01 and S02. (b) Leave-one-soil-out predictions against observations (RMSE 0.058). (c) LOSO RMSE of each predictor of $\phi_D$ (asterisk: eight soils; the other predictors are shown for all ten soils, and like-for-like eight-soil values are given in \cref{tab:loso}).}
\label{fig:mech}
\end{figure}

\begin{table}[H]\centering\small
\caption{Leave-one-soil-out (LOSO) predictive error of $\phi_D$ by predictor.}
\label{tab:loso}
\begin{tabular}{lrrl}
\toprule
Predictor & Soils & LOSO RMSE & Model\\\midrule
Mean only & 10 & 0.318 & constant\\
Mean only & 8 & 0.338 & constant\\
$\TAW$ & 8 & 0.088 & linear, clipped\\
$\Se(100)$ & 8 & 0.064 & linear, clipped\\
$n$ & 8 & 0.100 & linear, clipped\\
$\TAW$ & 10 & 0.183 & linear, clipped\\
$\Se(100)$ & 10 & 0.136 & linear, clipped\\
$n$ & 10 & 0.097 & linear, clipped\\
$\log_{10}K_s$ (linear) & 8 & 0.062 & linear, clipped\\
$K_s/(K_s+q^\ast)$ & 8 & 0.058 & \cref{eq:rate}, $q^\ast=71.5$\\\bottomrule
\end{tabular}
\end{table}

\subsection{Memory fate predicts the loss fate in independent experiments}
\label{sec:transfer}
The key test of transferability is whether $\phi_D$, measured with antecedent moisture, predicts $\psi_D$, measured in experiments that never varied antecedent moisture. It does: $\phi_D$ against $\psi_D^{E16}$ gives $\rho=0.85$ ($p=0.002$) and against $\psi_D^{E7}$ gives $\rho=0.89$ ($p=0.0014$; \cref{fig:rep}). The two loss-fate experiments agree with each other less well ($\rho=0.71$). The van~Genuchten $n$ predicts $\psi_D^{E7}$ well ($\rho=0.89$, $p=0.001$) but $\psi_D^{E16}$ only moderately ($\rho=0.70$, $p=0.03$, bootstrap CI $-0.06$ to 1.00), and the bootstrap interval includes zero. \cref{fig:robust} shows the permutation distribution for $n$ versus $\phi_D$ and the bootstrap intervals of all five associations. \emph{An important qualification}: $n$ predicts $\psi_D$ as well as $\phi_D$ does, so the data do not show that the memory fate is a better predictor of recoverability than the soil's own hydraulic shape parameter. What $\phi_D$ adds is an operational definition and a demonstration that the single mechanism (drainage competition) shows up under three different perturbations.

\begin{figure}[H]\centering
\includegraphics[width=\textwidth]{\figdir/p1/fig08_replication.pdf}
\caption{Cross-experiment transfer. (a) $\psi_D$ from E16 against $\phi_D$ from E9; (b) $\psi_D$ from the mulch lattice E7 against $\phi_D$; (c) E7 against E16. Each point is a soil, coloured by $n$ (as in \cref{fig:predictors}).}
\label{fig:rep}
\end{figure}

\begin{figure}[H]\centering
\includegraphics[width=\textwidth]{\figdir/p1/fig10_robustness.pdf}
\caption{Robustness of the key associations. (a) Permutation distribution of Spearman $\rho$ between $n$ and $\phi_D$ (solid line: observed; dotted: its mirror). (b) Bootstrap 95\,\% intervals of five rank associations (10\,000 resamples).}
\label{fig:robust}
\end{figure}

\begin{table}[H]\centering\small
\caption{Rank associations among soil-level diagnostics ($n=10$ soils). CI: percentile bootstrap.}
\label{tab:rank}
\begin{tabular}{lrrr}
\toprule
Pair & Spearman $\rho$ & Permutation $p$ & Bootstrap 95\,\% CI\\\midrule
$n$ vs.\ $\phi_D$ & 0.95 & 0.00015 & 0.69--1.00\\
$\phi_D$ vs.\ $\psi_D$ (E16) & 0.85 & 0.002 & 0.37--1.00\\
$\phi_D$ vs.\ $\psi_D$ (E7) & 0.89 & 0.0014 & 0.55--1.00\\
$n$ vs.\ $\psi_D$ (E16) & 0.70 & 0.032 & $-0.06$--1.00\\
$n$ vs.\ $\psi_D$ (E7) & 0.89 & 0.001 & 0.47--1.00\\\bottomrule
\end{tabular}
\end{table}

\subsection{Climate and management modify the recoverable share}
\label{sec:mod}
Because $\psi_D$ was computed at each of the 32 combinations of the other factors in E16, it can be recomputed at each level of a second factor, and the paired change across soils tested (\cref{fig:mod}, \cref{tab:mod}). Warming from $+2$ to $+6$\,\textdegree C lowers $\psi_D$ in all ten soils (mean change $-0.16$; sign test $p=0.002$), so a larger share of the extra evaporation is offset by irrigation or transpiration rather than drainage, and the recoverable share rises. In S01 $\psi_D$ falls from 0.41 to 0.17. Raising rainfall from 0.5 to 1.5 times the baseline raises $\psi_D$ in nine of nine soils with data in both levels (mean change $+0.42$; $p=0.004$). Wetter climates leave more water to drain, so the same evaporation saving is more likely to be a drainage saving. Trigger suction (+0.14, 7/10 soils up, $p=0.16$), dose ($-0.04$, 2/10 up, $p=0.32$) and nitrogen rate ($-0.008$, $p=0.064$, a negative control that should not influence the water partition) are not significant at $\alpha=0.05$.

\begin{figure}[H]\centering
\includegraphics[width=\textwidth]{\figdir/p1/fig07_modifiers.pdf}
\caption{Modifiers of the drainage share $\psi_D$ in E16. Each line is a soil (colour: $n$); (a) rainfall scale, (b) warming, (c) trigger suction, (d) irrigation dose.}
\label{fig:mod}
\end{figure}

\begin{table}[H]\centering\small
\caption{Paired effects of E16 factors on the drainage share $\psi_D$ of an evaporation change. ``Up'' counts soils in which $\psi_D$ is higher at the second level. $p$: exact two-sided sign test across soils.}
\label{tab:mod}
\begin{tabular}{lcrrr}
\toprule
Factor (low $\to$ high) & Mean change in $\psi_D$ & Soils up / n & $p$ & Verdict\\\midrule
Rain scale $0.5\to1.5$ & $+0.419$ & 9/9 & 0.004 & supported\\
Warming $+2\to+6$\,\textdegree C & $-0.161$ & 0/10 & 0.002 & supported\\
Trigger $50\to2000$\,cm & $+0.136$ & 7/10 & 0.160 & not significant\\
Dose $1.2\to30$\,mm & $-0.040$ & 2/10 & 0.322 & not significant\\
N rate $25\to150$ & $-0.008$ & 2/10 & 0.064 & negative control\\\bottomrule
\end{tabular}
\end{table}

\label{sec:e11}
The E16 contrasts use only two levels per factor. Experiment E11 provides a finer test for coarse loam (S01): a $10\times10\times10$ lattice of warming $\Delta T\in[-1,3.5]$\,\textdegree C, evaporation factor $E_p\in[0.3,1.2]$ and rain scale $P\in[0.5,1.4]$. For each $(\Delta T,P)$ pair, drainage is regressed on $\Ent$ across the ten $E_p$ levels, giving $\psi_D=-\mathrm{d}D/\mathrm{d}\Ent$ (\cref{fig:e11}). The surface is smooth and monotone: $\psi_D$ rises with rainfall (Spearman $\rho=0.88$ with $P$ over the 100 cells) and falls with warming ($\rho=-0.40$), ranging from 0.01 in the driest, warmest corner to 0.61 at $P=1.4$ and $\Delta T=-1$\,\textdegree C. Averaged over $\Delta T$, the drainage share rises from 0.013 at $P=0.5$ to 0.367 at $P=1.4$; averaged over $P$, it falls from 0.283 at $\Delta T=-1$\,\textdegree C to 0.036 at $+3.5$\,\textdegree C. This replicates, within a single soil and over a continuous range, the signs found across soils in E16 (\cref{tab:mod}). It also shows that the single number $\psi_D=0.295$ reported for S01 in \cref{tab:fates} is a point on a surface that spans a factor of about 50, so any soil-level value of the recoverable fraction should be quoted together with the climate for which it was computed. The magnitude of the warming effect depends on $P$: at $P=0.5$ the drainage share is near zero everywhere (0.01--0.02), so the recoverable share is close to one regardless of temperature, whereas at $P=1.4$ it falls from 0.61 to 0.09 across the warming range.

\begin{figure}[H]\centering
\includegraphics[width=\textwidth]{\figdir/p1/fig11_e11_surface.pdf}
\caption{Climate surface of the drainage share $\psi_D$ for coarse loam (S01) from the E11 lattice. (a) $\psi_D$ over rain scale and warming (white contours), (b) $\psi_D$ against $P$ for selected warming levels, (c) $\psi_D$ against warming for selected rain scales.}
\label{fig:e11}
\end{figure}

\subsection{Alternative formulations and what they establish}
\label{sec:fail}
Four alternative formulations were tested against the ensemble, and each outcome sharpens the definition of the memory fate. (A1) A universal curve of $D/W$ against $W/\TAW$ that collapses all soils does not exist: the curves remain ordered by soil, which establishes that recoverability must be defined per soil and is the reason $\phi_D$ is a soil-level diagnostic. (A2) Nitrogen leaching does not track drainage with a soil-independent coefficient; it depends on soil and fertiliser timing, so the water diagnostics are the transferable quantities and nitrogen is a soil-conditional downstream response. (A3) Trigger suction and irrigation dose are not first-order controls of recoverability (neither is significant), so the recoverable share is set by soil and climate and is robust to scheduling detail. (A4) Literature conductivities are available for eight of the ten soils, so the conductivity law is calibrated and validated on those eight, with the pore-size parameter $n$ providing the ten-soil comparison. The pore-size parameter $n$ and the memory fate carry comparable information on the loss fate; this is expected, because both are expressions of one pore structure, and \cref{sec:disc} sets out what the flux-based definition adds.

\subsection{Soil-group profiles of the three fates}
\paragraph{Sand (S03).} $\phi_D=0.85$, the highest of the set. With $\TAW=60$\,mm and $\mathrm{cap}_{FC}=4$\,mm the profile cannot hold antecedent water: release saturates at $h_0\approx100$\,cm (\cref{fig:memory}a). Both loss fates are extreme ($\psi_D^{E16}=1.21$, $\psi_I^{E16}<0$), and the mulch lattice returns $\psi_D^{E7}=0.71$ with $\psi_I^{E7}=0$, that is, none of the mulch saving appears as reduced irrigation. In this soil an evaporation saving is effectively an increase in deep drainage.
\paragraph{Loamy sand and sandy loam (S04, S05).} $\phi_D=0.74$ and 0.64. $\psi_D^{E7}=0.69$ and 0.66 with slightly negative $\psi_I^{E7}$. These two soils lie between sand and the loams in $n$ (2.28, 1.89) and also in behaviour; they are the soils for which the $K_s$ law has its largest absolute residuals after sand (\cref{fig:mech}b).
\paragraph{Intermediate soils (S06, S08, S02).} $\phi_D\approx0.26$--0.29, with $\psi_D^{E16}=0.47$, 0.53, 0.19. The coarse sand S02 is anomalous in two ways: it has $n=1.60$, close to the loams, and its drainage share is therefore similar ($\phi_D=0.265$), yet its $\psi_D^{E16}$ is the smallest of the set while its $\psi_D^{E7}$ (0.61) is among the largest. This is attributed to the reduced E16 design for S02 (37 valid combinations), and the S02 loss fates are read as indicative.
\paragraph{Fine and structured soils (S07, S01, S09, S10).} $\phi_D=0.02$--0.15; $\psi_D^{E16}\approx0.30$--0.36; $\rho_R\approx0.64$--0.71. These are the soils in which evaporation savings are most recoverable and in which irrigation compensates most strongly ($\psi_I^{E16}=0.53$--0.84).

\subsection{Sensitivity of the conclusions to analytical choices}
Three analytical choices could in principle drive the results. (i) \emph{Pooling within $P$ in \cref{eq:phi}} is essential, not cosmetic. A single regression of drainage on $\bar R$ over all 100 cells, ignoring $P$, gives values that differ from the pooled $\phi_D$ by 0.003--0.04 in most soils but by $-0.23$, $-0.48$ and $-0.93$ in S05, S04 and S03 (the rank correlation between the two estimates is only 0.43), because in free-draining soils rainfall changes both $\bar R$ and drainage strongly. The memory fate is therefore defined as the response to antecedent state \emph{at fixed rainfall}, and the results above should not be compared with unpooled slopes. (ii) \emph{Ratio-of-sums in \cref{eq:psi}}: a ratio of means is more stable than a mean of ratios because individual paired $\Delta\Ent$ can be small when $P$ is low; the bootstrap intervals in \cref{fig:fates}b reflect this. (iii) \emph{Soil-level correlations treat each soil as independent}: this inflates significance if soils are correlated through texture, which is why both permutation $p$ values and bootstrap intervals are reported and the interval is emphasised for $n$ versus $\psi_D^{E16}$, whose lower bound is below zero.

\begin{table}[H]\centering\small
\caption{Hypothesis register. the test and outcome of every formulation examined. Verdicts refer to the ten-soil simulated ensemble.}
\label{tab:hyp}
\begin{tabularx}{\textwidth}{lXlX}
\toprule
ID & Hypothesis & Verdict & Evidence\\\midrule
H0 & $\phi_D$ differs among soils and is ordered by $n$ & supported & $\rho=0.95$; \cref{tab:rank}\\
H1 & $\phi_D$ transfers to $\psi_D$ (E16, E7) & supported & $\rho=0.85$, 0.89\\
H2 & One-parameter $K_s$ law beats structural predictors out of sample & supported (8 soils) & LOSO RMSE 0.058 vs.\ 0.097--0.183\\
H3 & $\phi_D$ predicts recoverability better than $n$ itself & comparable to $n$ & $n$ vs.\ $\psi_D^{E7}$: $\rho=0.89$\\
H4 & Warming raises recoverable share & supported & 10/10 soils; E11 surface\\
H5 & Wetter climate lowers recoverable share & supported & 9/9 soils; E11 surface\\
H6 & Trigger or dose is a first-order modifier & not first-order & $p=0.16$, 0.32\\
H7 & Nitrogen has no effect (negative control) & as expected & $-0.008$\\
H8 & Universal $D/W$--$W/\TAW$ collapse & soil-specific & curves remain soil-ordered\\
H9 & $N$ leaching tracks drainage with a common coefficient & soil-specific & soil-dependent\\\bottomrule
\end{tabularx}
\end{table}

%==========================================================
\section{Results of the accounting audit}
\label{sec:resultsB}
\subsection{The pipeline recovers the water balance (positive control)}
The singular spectrum of the seven-flux matrix falls steeply, with a smallest singular value of 5.5 compared with about 33 under column permutation (\cref{fig:invariants}a), and the corresponding right singular vector places weight on rain, irrigation, ET and drainage with opposite signs for inputs and outputs (\cref{fig:invariants}b). For the four-flux subsystem the null direction, scaled to $c_{\mathit{ET}}=-1$, is
\begin{equation}
\mathbf c=(c_{P_r},c_I,c_{\mathit{ET}},c_D)=(1.037,\;0.957,\;-1,\;-1.036),
\label{eq:law}
\end{equation}
with bootstrap 95\,\% intervals 1.031--1.045 for rain, 0.940--0.974 for irrigation and $-1.052$ to $-1.023$ for drainage, and a smallest singular value of 5.9 against 33.3 under permutation. The pipeline therefore recovers the conservation law $P_r+I-\mathit{ET}-D\approx0$ with every coefficient within 5\,\% of the exact value, and the permutation null shows that the null direction is not an artefact of scale. \emph{However, the bootstrap intervals do not contain the exact values}: rain $\pm$ 1 lies below the rain interval, and irrigation lies above its interval. The deviation reflects the nonzero storage term and other unreported fluxes: excluding E8 and E9, whose initial moisture is deliberately unbalanced, gives $(1.021,\,0.969,\,-1,\,-1.019)$ with a smallest singular value of 3.0, closer to the exact balance but still not exact.

\begin{figure}[H]\centering
\includegraphics[width=\textwidth]{\figdir/p2/fig02_invariants.pdf}
\caption{Linear invariants. (a) Singular values of the seven-flux matrix for the real data and for column-permuted data. (b) Coefficients of the smallest right singular vector (scaled so $c_{\mathit{ET}}=-1$) with bootstrap intervals. (c) Storage residual $\Delta S/W$ by experiment: E8 and E9 do not close. (d) Unattributed share of ET, $R/\mathit{ET}$, by experiment.}
\label{fig:invariants}
\end{figure}

The result illustrates both the promise and the discipline needed for law discovery. The method finds the law that is there, without being told which variables are inputs and which are outputs; but its coefficients should not be over-interpreted, because the third decimal of a coefficient reflects what the data do \emph{not} close, not what the law is.

Fitting the four-flux null vector separately within each experiment (\cref{tab:svdexp}) tests whether the pooled vector \cref{eq:law} hides heterogeneity. In 11 of the 14 experiments with irrigation all coefficients are within about 5\,\% of the exact values $(1,1,-1,-1)$ (E11 is at 5.1\,\% for drainage; E16, with the extreme corners of the factorial, gives 0.901 for irrigation and $-0.919$ for drainage, E6 gives 0.961, 0.927 and $-0.869$; and E12 gives drainage $-1.116$ are the exceptions). The residual scale $\sigma_q/\sqrt N$ is 0.001--0.017 in the well-closing experiments and 0.081 in E16, 0.030 in E6, 0.166 in E8 and 0.219 in E9. E0 (no irrigation variation within the 10 baseline runs), E8 and E9 (rainfed) have a constant zero irrigation column and the invariant is formed from three fluxes; for E8 the drainage coefficient is $+15.5$, a sign of an ill-conditioned decomposition in an experiment in which storage dominates the balance, which is a cautionary example of how an unconstrained search can return a numerically valid but physically meaningless vector. The method therefore identifies its own failure by the magnitude of $\sigma_q/\sqrt N$ and by physical-plausibility checks on the coefficients, which is recommended as a rule.

\begin{table}[H]\centering\small
\caption{Four-flux null vector fitted within each experiment (scaled to $c_{\mathit{ET}}=-1$); ``--'' marks a flux with no variation in the experiment (irrigation in E0, E8, E9). $\hat\sigma=\sigma_q/\sqrt N$ is the residual scale in units of the column root-mean-square.}
\label{tab:svdexp}
\begin{tabular}{lrrrrr}
\toprule
Experiment & Runs & $c_{P_r}$ & $c_{I}$ & $c_{D}$ & $\hat\sigma$\\\midrule
E0 & 10 & 1.025 & -- & $-$1.041 & 0.004\\
E1 & 189 & 1.009 & 0.969 & $-$0.979 & 0.017\\
E2 & 197 & 1.022 & 0.983 & $-$1.027 & 0.004\\
E3 & 200 & 1.016 & 0.995 & $-$1.012 & 0.004\\
E4 & 96 & 1.006 & 0.989 & $-$0.983 & 0.003\\
E5 & 300 & 1.016 & 0.993 & $-$1.027 & 0.002\\
E6 & 193 & 0.961 & 0.927 & $-$0.869 & 0.030\\
E7 & 200 & 1.008 & 0.992 & $-$0.990 & 0.003\\
E8 & 1000 & 1.241 & -- & 15.496 & 0.166\\
E9 & 987 & 1.182 & -- & $-$1.249 & 0.219\\
E10 & 954 & 1.026 & 0.972 & $-$1.038 & 0.010\\
E11 & 1000 & 0.998 & 1.024 & $-$1.051 & 0.004\\
E12 & 1000 & 1.016 & 1.005 & $-$1.116 & 0.005\\
E13 & 975 & 1.016 & 0.984 & $-$1.008 & 0.004\\
E14 & 1000 & 0.982 & 1.035 & $-$0.974 & 0.001\\
E15 & 1000 & 1.009 & 0.966 & $-$0.971 & 0.003\\
E16 & 740 & 1.018 & 0.901 & $-$0.919 & 0.081\\
\bottomrule
\end{tabular}
\end{table}

The closure residual $\varepsilon$ (\cref{eq:closure}) has a median absolute value of 0.9\,\% of water input over all valid runs and 0.65\,\% excluding E8 and E9 (95th percentile 2.3\,\%), so for 15 of the 17 experiments the mass balance holds to within a few percent. In E8 and E9, which vary initial moisture, the median absolute residual is 16\,\% of $W$, with the residual reaching $\pm0.3$--0.6 in the tails (\cref{fig:invariants}c). This is expected, because these experiments deliberately impose a non-equilibrium initial storage that is released over the 20-year window, but it matters for discovery: a conservation law fitted without storage will appear to be violated in exactly the experiments where hydraulic memory is strongest. A global law should therefore include storage explicitly or be fitted only on near-equilibrium experiments, and the two cases give different coefficients (\cref{eq:law} and the text following it).

\subsection{A unattributed term in the ET partition}
The ET components that Daisy reports do not add up to ET. The remainder $R$ (\cref{eq:R}) averages 164\,mm\,yr$^{-1}$, with an interquartile range of 103--223\,mm\,yr$^{-1}$, and equals 30\,\% of ET (median 29\,\%; range 2--88\,\%). Its share varies between experiments from 14\,\% (E6, E8) to 41\,\% (E7, E14; \cref{fig:invariants}d, \cref{tab:R}). The dependence of $R$ on management is informative (\cref{fig:hidden}). In E4 (methods), $R$ is 105\,mm\,yr$^{-1}$ for every subsurface-drip configuration (5--60\,cm), 183 for surface and 166 for overhead irrigation, so that \emph{where} water is applied matters while the drip depth does not (\cref{fig:hidden}a). In the rainfed E9 lattice, $R=50.9+0.121\,P_r$ (\cref{fig:hidden}b), i.e.\ about 12\,\% of rainfall, and in E11 (coarse loam)
\begin{equation}
\ln R=5.75-0.355\ln P+0.863\ln E_p+0.136\,\Delta T,\qquad R^2=0.979,
\label{eq:Rlaw}
\end{equation}
where $P$ is the rain scale, $E_p$ the soil-evaporation factor and $\Delta T$ the temperature shift, so that $R$ scales almost linearly with the soil-evaporation factor, decreases with rain scale (elasticity $-0.36$) and increases with warming. The mulch experiments E7 and E14 have the highest median shares (41\,\%) and the cross-experiment sparse model for $\ln R$ selects $\ln(1+I)$, $n$, the mulch capacity and $\Delta T$.

\begin{figure}[H]\centering
\includegraphics[width=\textwidth]{\figdir/p2/fig03_hidden_term.pdf}
\caption{The unattributed term $R$. (a) $R$ and canopy evaporation in E4 by irrigation method. (b) $R$ against annual rain in E9 (rainfed, light points), with the linear fit and the ten E0 baselines.}
\label{fig:hidden}
\end{figure}

\begin{table}[H]\centering\small
\caption{Unattributed term $R$ by experiment. median $R$ (mm\,yr$^{-1}$), median share of ET, and median storage residual $\Delta S/W$.}
\label{tab:R}
\begin{tabular}{lrrrr}
\toprule
Experiment & Runs & $R$ (mm) & $R/\mathit{ET}$ & $\Delta S/W$\\\midrule
E0 Baselines & 10 & 107 & 0.24 & $-0.021$\\
E1 Doses & 189 & 118 & 0.21 & $-0.001$\\
E2 Trigger pressure & 197 & 168 & 0.30 & $-0.009$\\
E3 Sensor depth & 200 & 202 & 0.32 & $-0.009$\\
E4 Methods & 96 & 103 & 0.19 & $-0.003$\\
E5 SHP dynamics & 300 & 162 & 0.29 & $-0.007$\\
E6 Mediterranean & 193 & 70 & 0.14 & $-0.012$\\
E7 Mulching & 200 & 226 & 0.41 & $-0.006$\\
E8 Rain intensity$\times$moisture & 1000 & 39 & 0.14 & $-0.133$\\
E9 Soil moisture$\times$rain & 987 & 81 & 0.37 & $-0.014$\\
E10 Amount, timing, soil & 954 & 118 & 0.23 & $-0.011$\\
E11 Temperature$\times E_p\times$rain & 1000 & 296 & 0.38 & $-0.009$\\
E12 Sensor, trigger, rain & 1000 & 218 & 0.32 & $-0.008$\\
E13 N$\times$water & 975 & 125 & 0.28 & $-0.006$\\
E14 Mulch$\times$rain & 1000 & 262 & 0.41 & $0.001$\\
E15 Method, depth, amount & 1000 & 103 & 0.17 & $0.003$\\
E16 Edge cases & 740 & 157 & 0.30 & $-0.010$\\\bottomrule
\end{tabular}
\end{table}

What is $R$? The data do not identify it, but they constrain candidates. A process that is independent of drip depth, larger under surface and overhead application, proportional to the soil-evaporation factor, increased under mulch and weakly increasing with warming is consistent with evaporation from the \emph{surface layer} (mulch or litter and ponded water) that is not reported in the soil-evaporation or canopy-evaporation columns. That reading is a hypothesis; the model code has not been inspected to confirm which Daisy state variables sum into the model's ET output, and a software-level audit of the output definitions would be required to turn the hypothesis into a statement. The practical consequence is nonetheless clear: any analysis that treats $T+E_s+E_c$ as ET (as many reduced-form ET-partitioning analyses do) omits about 30\,\% of the flux, and the partition fractions reported for the ensemble depend on how $R$ is allocated.

A sparse model for $R$ across all experiments has limited out-of-sample skill. The forward-selected model $\ln R=3.35+0.194\ln(1+I)+0.492\,n+0.145\,\mathrm{cap}+0.238\,\Delta T$ has an in-sample $R^2=0.61$ on the log scale but a leave-one-experiment-out error of 0.51 in $\ln R$, equal to a level-scale $R^2$ of only 0.30 (RMSE 81\,mm) versus 0.50 (RMSE 69\,mm) for gradient boosting. Within a single experiment, \cref{eq:Rlaw} is excellent; across experiments it is not, which implies that $R$ depends on interactions among irrigation method, mulch state and demand that the library does not contain.

\subsection{A dominant drainage axis, with its variance share quantified by bootstrap}
The log-ratio transform requires strictly positive components, and drainage is zero or negative (net capillary rise) in 1\,429 of the 10\,041 runs (14\,\%) and below 1\,mm\,yr$^{-1}$ in 1\,947. The treatment of these zeros is therefore not a detail. The first analysis replaced all components below $10^{-3}$\,mm\,yr$^{-1}$ by that floor, which gave explained variances of 88.6\,\%, 7.3\,\% and 3.3\,\% for the first three components. Because $\ln(10^{-3})$ lies far from the logs of the other components, the runs with zero drainage then dominate the first axis. Repeating the analysis with floors of 0.1, 1, 5 and 10\,mm\,yr$^{-1}$ gives a first-component share of 74.3, 62.7, 56.5 and 52.3\,\% (\cref{fig:manifold}a, \cref{tab:clipsens}). The result that a single drainage axis dominates is robust, since the loading of drainage on PC1 stays between $+0.82$ and $+0.88$ for every floor and also when only runs with drainage above 5 or 20\,mm are used (PC1 share 57\,\% and 52\,\%), but its \emph{magnitude} is a function of the data-handling choice and the figure of 88.6\,\% should not be quoted as a property of the system. The number of components needed to explain 90\,\% of the variance is two to three in all cases (three components carry 92--98\,\% for floors from $10^{-3}$ to 10\,mm, with 97.3\,\% at the 1\,mm floor adopted for the loadings, \cref{fig:manifold}b,c).

With the 1\,mm floor, the first component (62.7\,\%) is a drainage-versus-evaporation axis: drainage loads $+0.84$, soil evaporation $-0.37$, unattributed $R$ $-0.34$, transpiration $-0.20$ and canopy evaporation $+0.06$. Its scores correlate with the drainage fraction $D/W$ at Spearman $\rho=0.98$, with $\ln(W/\TAW)$ at 0.48 and with $n$ at 0.45. The second component (24.8\,\%) separates canopy evaporation ($+0.77$) from soil evaporation ($-0.56$), and the third (9.8\,\%) opposes $R$ ($+0.83$) to soil evaporation ($-0.43$) and transpiration ($-0.26$). The choice of the 1\,mm floor was made \emph{after} seeing the sensitivity and is justified only as a value close to the resolution of annual flux accounting; the full sensitivity is reported so that readers can judge. The practical conclusion is that the drainage fraction is, to a first approximation, the leading summary of water partitioning (its score correlation is 0.98), and that a second axis, the canopy-versus-soil evaporation split, is needed for a complete description.

\begin{figure}[H]\centering
\includegraphics[width=\textwidth]{\figdir/p2/fig04_manifold.pdf}
\caption{Partition manifold and its sensitivity to zero handling. (a) Variance explained by the first three clr-PCA components for five zero-replacement floors. (b) Scores on the first two components (1\,mm floor) coloured by $\ln(W/\TAW)$. (c) Loadings of the five fates on the first two components.}
\label{fig:manifold}
\end{figure}

\begin{table}[H]\centering\small
\caption{Sensitivity of the clr-PCA to the zero-replacement floor (variance explained, \%).}
\label{tab:clipsens}
\begin{tabular}{lrrrr}
\toprule
Floor (mm\,yr$^{-1}$) & PC1 & PC2 & PC3 & PC4\\\midrule
$10^{-3}$ & 88.6 & 7.3 & 3.3 & 0.8\\
0.1 & 74.3 & 17.3 & 6.6 & 1.7\\
1 & 62.7 & 24.8 & 9.8 & 2.6\\
5 & 56.5 & 25.1 & 14.8 & 3.6\\
10 & 52.3 & 27.1 & 16.8 & 3.8\\\bottomrule
\end{tabular}
\end{table}

Averaged over the 15 near-equilibrium experiments, the five fates (\cref{fig:fatesoil}a) show the expected gradient: drainage takes 34.5\,\% of water input in sand (S03), 28.5\,\% in loamy sand and 5.4\,\% in coarse loam (S01), 8.5\,\% in silty clay loam, while transpiration rises from 19\,\% in sand to 31--33\,\% in the loams. The unattributed term, in contrast, is remarkably constant across soils: $R$ takes 22--26\,\% of $W$ in nine of ten soils (31\,\% in coarse loam, which has the mulch and temperature lattices) and its share of ET is 23--39\,\% (\cref{fig:fatesoil}c), the largest in sand, where there is least else to evaporate. Since $R$ is almost soil-independent in share of $W$, it behaves as a surface-driven term rather than a soil-hydraulic one, consistent with the surface-layer hypothesis. The stacked fractions sum to 0.96--1.01 of $W$ (the remainder being storage change and runoff). Panel~b is a reminder that coarse loam contributes 4\,419 of the near-equilibrium runs while each other soil contributes 290--419.

\begin{figure}[H]\centering
\includegraphics[width=\textwidth]{\figdir/p2/fig07_fates_by_soil.pdf}
\caption{Water fates by soil in the 15 near-equilibrium experiments (E8, E9 excluded). (a) Mean fate shares of water input $W$, soils ordered by decreasing $n$. (b) Number of valid runs per soil. (c) Median share of the unattributed term in ET.}
\label{fig:fatesoil}
\end{figure}

\subsection{A compact evapotranspiration law with out-of-sample skill}
Forward selection over the library selects three terms for ET (\cref{tab:laws}): $\ln W$, $n$ and $\ln(1+I)$, giving
\begin{equation}
\ln\mathit{ET}=2.503+0.646\,\ln W-0.299\,n+0.029\,\ln(1+I),
\label{eq:ETlaw}
\end{equation}
with ordinary-least-squares standard errors $0.027$, $0.0051$, $0.0045$ and $0.0011$ (which understate uncertainty because runs are not independent) and in-sample $R^2=0.908$. Leave-one-experiment-out $R^2=0.898$ and leave-one-soil-out $R^2=0.904$, against 0.876--0.879 (LOEO) and 0.881--0.894 (LOSO) for gradient boosting in two independent runs, so the three-term law matches or exceeds the benchmark under both groupings (\cref{fig:ETlaw}a,b). The RMSE in leave-one-experiment-out prediction is 68.8\,mm\,yr$^{-1}$ (MAPE 9.2\,\%). The fitted coefficients are stable across the 17 folds (\cref{fig:ETlaw}c; coefficient means 2.51, 0.645, $-0.298$, 0.029 with standard deviations of 0.077, 0.017, 0.017, 0.0023).

Equation~\eqref{eq:ETlaw} has a straightforward physical reading. The exponent on $W$ (0.646) means that ET rises sublinearly with water supply, the signature of an energy- or storage-limited regime also encoded in Budyko-type curves \citep{Budyko1974,Milly1994,Zhang2001}: a 10\,\% increase in supply raises ET by 6.5\,\%, and the remainder goes to drainage, runoff and storage. The coefficient on $n$ (-0.299) implies a factor $\exp(-0.299\times1.45)=0.65$ in ET between the most uniform (S03, $n=2.68$) and the most graded (S10, $n=1.23$) soil at the same $W$: soils that drain freely retain less water for evaporation, in agreement with the hydraulic-memory results of the memory analysis (\cref{sec:results}). The small positive coefficient on $\ln(1+I)$ (0.029) indicates that, at equal total input, a larger irrigated share raises ET slightly, plausibly through the wetting of the surface that increases $R$ and $E_s$. The law is \emph{not} a fundamental conservation law, nor a model of mechanism; it is a compact summary of how the model's ET responds across 17 designs.

\begin{figure}[H]\centering
\includegraphics[width=\textwidth]{\figdir/p2/fig05_ET_law.pdf}
\caption{Compact ET law. (a) Leave-one-experiment-out predictions of the three-term law versus Daisy ET. (b) Gradient-boosting benchmark. (c) Fitted coefficients across the 17 LOEO folds. (d) LOEO MAPE by experiment. (e) The law at $I=200$\,mm for four soils (grey dashed: ET $=W$). (f) Drainage derived from closure, $W-\widehat{\mathit{ET}}$, versus Daisy drainage ($R^2=0.69$, RMSE 64\,mm; E8 and E9 excluded).}
\label{fig:ETlaw}
\end{figure}

The 17 experiments were resampled with replacement (2\,000 cluster-bootstrap resamples) and the cross-validated predictions of the law were compared with those of a gradient-boosting model that uses the same three inputs ($\ln W$, $n$, $\ln(1+I)$; LOEO $R^2=0.866$). The paired difference in $R^2$ is $+0.032$ with a 95\,\% interval of 0.004--0.087, so the law is better than a flexible learner given the same information in this resampling, with the caveat that the law's terms were selected with the same folds. With the \emph{full} library the benchmark reaches 0.876--0.879, and the two are regarded as comparable rather than claiming superiority.

The error is not uniform (\cref{tab:byexp}). Mean absolute percentage errors are 2.6--5.8\,\% in most one- and two-factor experiments, but 16.6\,\% in E16 (edge cases: extreme combinations), 16.7\,\% in E8 and 27.4\,\% in E9, the two experiments that violate storage closure, and 11.7\,\% in E5 (time-varying soil hydraulic properties, where a single static $n$ is wrong by construction). The law is therefore valid where the water balance is near equilibrium and soil hydraulics are static; it should not be extrapolated to initial-state perturbations or to dynamic soil structure.

\begin{table}[H]\centering\small
\caption{Leave-one-experiment-out performance of \cref{eq:ETlaw} by experiment: MAPE (\%) and mean error (predicted minus Daisy, mm\,yr$^{-1}$).}
\label{tab:byexp}
\begin{tabular}{lrrr}
\toprule
Experiment & Runs & MAPE (\%) & Mean error (mm)\\\midrule
E0 Baselines & 10 & 5.8 & $-25.2$\\
E1 Doses & 189 & 3.8 & 9.8\\
E2 Trigger pressure & 197 & 3.0 & $-6.6$\\
E3 Sensor depth & 200 & 2.6 & $-8.0$\\
E4 Methods & 96 & 3.6 & 9.4\\
E5 SHP dynamics & 300 & 11.7 & $-63.8$\\
E6 Mediterranean & 193 & 8.5 & 41.5\\
E7 Mulching & 200 & 4.7 & 10.6\\
E8 Rain intensity & 1000 & 16.7 & $-11.3$\\
E9 Soil moisture & 987 & 27.4 & 12.9\\
E10 Amount, timing & 954 & 3.3 & 1.3\\
E11 Temperature, $E_p$, rain & 1000 & 9.0 & $-62.6$\\
E12 Sensor, trigger, rain & 1000 & 3.3 & $-15.2$\\
E13 N, water & 975 & 3.5 & $-3.7$\\
E14 Mulch, rain & 1000 & 4.0 & 7.5\\
E15 Method, depth, amount & 1000 & 5.0 & 29.0\\
E16 Edge cases & 740 & 16.6 & 34.5\\\bottomrule
\end{tabular}
\end{table}

\Cref{tab:closure2} gives the closure residual $\varepsilon$ for every experiment. In 15 experiments the median absolute residual is below 2.1\,\% of $W$ and the 95th percentile below 4\,\% (E6, the Mediterranean-pollution scenarios, is the exception at 10.7\,\%); in E8 and E9 the median absolute residual is about 16\,\% and the 95th percentile exceeds 100\,\% of $W$. For E8 the mean (signed) residual is negative, $-13\,\%$: the profile is gaining water, that is, the run begins drier than the equilibrium state and refills. The law \cref{eq:ETlaw} was fitted on water inputs from 124 to 2\,762\,mm\,yr$^{-1}$ (ET of 95--1\,600\,mm\,yr$^{-1}$), $n$ from 1.23 to 2.68, in a temperate climate, and it should be applied only within this domain and with near-equilibrium storage.

\begin{table}[H]\centering\small
\caption{Closure residual $\varepsilon=(W-\mathit{ET}-D)/W$ by experiment (per cent of water input): median signed, median absolute and 95th percentile of absolute residual.}
\label{tab:closure2}
\begin{tabular}{lrrrr}
\toprule
Experiment & Runs & Median $\varepsilon$ & Median $|\varepsilon|$ & 95th pct.\ $|\varepsilon|$\\\midrule
E0 & 10 & -2.06 & 2.06 & 2.9\\
E1 & 189 & -0.12 & 0.33 & 1.2\\
E2 & 197 & -0.87 & 0.87 & 2.3\\
E3 & 200 & -0.87 & 0.87 & 2.0\\
E4 & 96 & -0.33 & 0.34 & 1.2\\
E5 & 300 & -0.74 & 0.74 & 1.9\\
E6 & 193 & -1.19 & 1.57 & 10.7\\
E7 & 200 & -0.57 & 0.57 & 1.2\\
E8 & 1000 & -13.35 & 15.76 & 105.8\\
E9 & 987 & -1.40 & 16.87 & 108.5\\
E10 & 954 & -1.08 & 1.17 & 2.5\\
E11 & 1000 & -0.93 & 0.94 & 1.8\\
E12 & 1000 & -0.82 & 0.83 & 1.8\\
E13 & 975 & -0.64 & 0.64 & 2.3\\
E14 & 1000 & 0.09 & 0.40 & 1.5\\
E15 & 1000 & 0.34 & 0.34 & 0.4\\
E16 & 740 & -1.01 & 1.22 & 3.8\\
\bottomrule
\end{tabular}
\end{table}

The ensemble is heavily imbalanced: 5\,519 of the 10\,041 runs (55\,\%) are coarse loam (S01) because the large lattices E8, E11, E12, E14 and E15 use only that soil, so soil-level behaviour in the other nine soils rests on about 400--520 runs each. The LOEO residuals of \cref{eq:ETlaw} by soil (\cref{tab:soilres}) show that the law over-predicts ET in the fine soils (mean error $+22$, $+20$ and $+31$\,mm in S08, S09 and S10, where positive means over-prediction relative to Daisy, since errors are predicted minus Daisy) and has the largest percentage errors in sand (15.1\,\%) and loamy sand (13.7\,\%), the soils for which $n$ is most extreme. A single shape parameter therefore does not remove all soil dependence, and the MAPE of about 7.5--15\,\% by soil is consistent with the overall 9.2\,\%.

\begin{table}[H]\centering\small
\caption{LOEO residuals of the ET law by soil (predicted minus Daisy, mm\,yr$^{-1}$). MAPE: mean absolute percentage error; RMSE in mm\,yr$^{-1}$.}
\label{tab:soilres}
\begin{tabular}{llrrrr}
\toprule
Soil & Texture & Runs & Mean error & MAPE (\%) & RMSE\\\midrule
S01 & Coarse loam & 5519 & -10.7 & 7.5 & 58.9\\
S02 & Coarse sand & 390 & 15.8 & 8.6 & 69.8\\
S03 & Sand & 506 & -13.6 & 15.1 & 63.8\\
S04 & Loamy sand & 519 & -4.8 & 13.7 & 57.9\\
S05 & Sandy loam & 519 & -3.8 & 10.9 & 56.1\\
S06 & Loam & 519 & -0.2 & 9.9 & 77.5\\
S07 & Silt loam & 519 & 0.4 & 10.6 & 106.2\\
S08 & Sandy clay loam & 519 & 22.2 & 10.2 & 57.8\\
S09 & Clay loam & 519 & 19.7 & 9.9 & 92.8\\
S10 & Silty clay loam & 512 & 31.2 & 11.6 & 107.4\\
\bottomrule
\end{tabular}
\end{table}

\begin{table}[H]\centering\small
\caption{Forward-selection path for the ET law. LOEO $R^2_{\mathrm{cv}}$ after each added term; and for other targets (first terms).}
\label{tab:path}
\begin{tabular}{llr}
\toprule
Target & Term added & $R^2_{\mathrm{cv}}$ (LOEO)\\\midrule
$\ln\mathit{ET}$ & 1 ($\ln W$) & 0.801\\
 & 2 ($n$) & 0.895\\
 & 3 ($\ln(1+I)$) & 0.898\\
$\ln T$ & 1 ($\ln W$) & 0.581\\
 & 2 ($n$) & 0.641\\
$\ln E_s$ & 1 ($n$) & 0.070\\
 & 2 (mulch) & 0.237\\
 & 3 ($\ln P_r$) & 0.308\\
$\ln E_c$ & 1 ($\ln W$) & 0.641\\
 & 2 ($\ln P_r$) & 0.651\\
 & 3 ($I/W$) & 0.668\\
$\ln$ yield & 1 ($\ln W$) & 0.531\\
 & 2 ($\ln(1+I)$) & 0.549\\
 & 3 ($\ln P_r$) & 0.599\\
 & 4 (mulch) & 0.604\\
$\ln$ DPF & 1 ($\ln(W/\TAW)$) & 0.193\\
 & 2 ($\ln\alpha$) & 0.260\\
 & 3 ($I/W$) & 0.422\\
 & 4 ($\Delta T$) & 0.433\\\bottomrule
\end{tabular}
\end{table}

If E8 and E9, which break closure, are removed from both fitting and validation (15 experiments, 8\,054 runs), the same three terms give LOEO $R^2=0.882$, LOSO $R^2=0.881$, RMSE 65.9\,mm and MAPE 5.8\,\%, so the law is more accurate (relative error 5.8\,\% versus 9.2\,\%) in the domain for which it is intended. The marginal contribution of the terms is clearer in this subset: $\ln W$ alone gives LOEO $R^2=0.600$ (RMSE 138\,mm), adding $n$ raises it to 0.857 (RMSE 74\,mm) and adding $\ln(1+I)$ to 0.882. A gradient-boosting model restricted to the same three inputs reaches LOEO $R^2=0.778$, below the law, so the advantage of the compact form is not that it uses more information but that its functional form (power law in $W$, exponential in $n$) extrapolates across experiments better than piecewise-constant trees, which cannot extrapolate beyond the range of the training experiments. (This comparison uses three inputs and 15 folds and is not directly comparable with the full-library benchmark of \cref{tab:laws}.)

\Cref{tab:example} applies \cref{eq:ETlaw} to three combinations of water input and irrigation in three soils. At equal input, ET is about 35\,\% lower in sand than in silty clay loam ($263$ versus $406$\,mm at $W=400$\,mm), because water not held in the root zone drains; the implied drainage plus storage change $W-\mathit{ET}$ is 137\,mm in sand and negative-to-small in the fine soils, in line with the memory results. The example is illustrative; its uncertainty is the RMSE of 66--69\,mm.

\begin{table}[H]\centering\small
\caption{ET (mm\,yr$^{-1}$) predicted by \cref{eq:ETlaw} for three water inputs $W$ with irrigation $I$.}
\label{tab:example}
\begin{tabular}{lrrr}
\toprule
 & $W=400$, $I=0$ & $W=650$, $I=200$ & $W=1000$, $I=300$\\\midrule
S03 sand ($n=2.68$) & 263 & 420 & 561\\
S01 coarse loam ($n=1.353$) & 391 & 624 & 834\\
S10 silty clay loam ($n=1.23$) & 406 & 648 & 865\\\bottomrule
\end{tabular}
\end{table}

\label{sec:failures}
\Cref{tab:laws} collects all sparse laws that were fitted. Transpiration is moderately predictable by a two-term law ($\ln T=-1.329+1.127\ln W-0.501\,n$; LOEO $R^2=0.641$, equal to gradient boosting 0.642). Canopy evaporation is predicted at LOEO 0.668 (gradient boosting 0.691) by $\ln E_c=-13.3-0.470\ln W+3.157\ln P_r+4.240\,I/W$, whose large coefficients reflect the narrow range of $E_c$ and should not be interpreted mechanistically. Crop yield is predicted at LOEO 0.604 (gradient boosting 0.650). In contrast, soil evaporation is poorly predicted by sparse terms ($\ln E_s=3.271-0.770\,n-0.380\,\mathrm{mulch}+0.462\ln P_r$; LOEO 0.308 versus 0.547 for gradient boosting; LOSO 0.348 versus 0.476), indicating that interactions that the library does not contain control $E_s$. The drainage fraction (DPF $=D/W$), the natural target given the PCA result, is the worst: the best sparse model with exogenous predictors has a LOEO $R^2$ of 0.43 and a LOSO $R^2$ of 0.22, against 0.50 (LOEO) for gradient boosting, and allowing partition predictors in the library raises gradient boosting to 0.59 but does not change the sparse law. The shortfall is a negative result for compact drainage laws.

\begin{table}[H]\centering\small
\caption{Sparse laws and benchmark. cross-validated $R^2$ on the log scale. Selected terms in order of selection.}
\label{tab:laws}
\footnotesize\setlength{\tabcolsep}{4pt}
\begin{tabular}{llrrrr}
\toprule
Target & Selected terms & Sparse LOEO & Sparse LOSO & GBM LOEO & GBM LOSO\\\midrule
$\ln\mathit{ET}$ & $\ln W$, $n$, $\ln(1+I)$ & 0.898 & 0.904 & 0.876--0.879 & 0.881--0.894\\
$\ln T$ & $\ln W$, $n$ & 0.641 & 0.652 & 0.642 & 0.659\\
$\ln E_s$ & $n$, mulch, $\ln P_r$ & 0.308 & 0.348 & 0.547 & 0.476\\
$\ln E_c$ & $\ln W$, $\ln P_r$, $I/W$ & 0.668 & 0.669 & 0.691 & 0.806\\
$\ln$ yield & $\ln W$, $\ln(1+I)$, $\ln P_r$, mulch & 0.604 & 0.617 & 0.650 & 0.743\\
$\ln$ DPF & $\ln(W/\TAW)$, $\ln\alpha$, $I/W$, $\Delta T$ & 0.433 & 0.217 & 0.498 & --\\
$\ln R$ & $\ln(1+I)$, $n$, cap, $\Delta T$ & 0.41$^{\dagger}$ & 0.43$^{\dagger}$ & 0.43--0.47$^{\dagger}$ & --\\\bottomrule
\end{tabular}

\smallskip\noindent\footnotesize$^{\dagger}$ $R^2$ on the log scale for $R$; level-scale $R^2$ is 0.30 (sparse) and 0.50--0.51 (gradient boosting).
\end{table}

Because the balance holds to within a few percent in 15 of 17 experiments, drainage can be derived by closure as $\hat D=W-\widehat{\mathit{ET}}$, which converts a hard prediction problem (drainage) into one with a good compact law (ET). On the 8\,054 runs outside E8 and E9 this yields $R^2=0.69$ and RMSE 64\,mm against Daisy drainage (\cref{fig:ETlaw}f), and negative predicted drainage in 4.8\,\% of runs. The accuracy is limited by the relative size of the quantities and is not uniform (\cref{tab:dbins}): the mean absolute error is smallest, 21--28\,mm, for Daisy drainage between 10 and 200\,mm (relative error 19--66\,\%), but the closure-derived value overpredicts by 32\,mm on average for drainage below 10\,mm (1\,944 runs) and \emph{underpredicts} by 69\,mm and 308\,mm for drainage of 200--500 and above 500\,mm. The ET law is smooth in $W$ and cannot represent the large drainage events, so closure-derived drainage is a useful consistency check for moderate drainage and not a predictor of extremes.

\begin{table}[H]\centering\small
\caption{Closure-derived drainage $W-\widehat{\mathit{ET}}$ (LOEO ET law) against Daisy drainage by drainage class (E8, E9 excluded). Bias = predicted minus Daisy.}
\label{tab:dbins}
\begin{tabular}{lrrr}
\toprule
Daisy drainage (mm\,yr$^{-1}$) & Runs & Bias (mm) & Mean absolute error (mm)\\\midrule
$<10$ & 1944 & $+31.7$ & 39.7\\
10--50 & 2631 & $+8.2$ & 21.0\\
50--100 & 1736 & $-2.2$ & 25.6\\
100--200 & 1284 & $-13.7$ & 28.3\\
200--500 & 377 & $-68.8$ & 104.0\\
$>500$ & 82 & $-308.1$ & 318.6\\\bottomrule
\end{tabular}
\end{table}

\subsection{Irrigation compensation is regime-specific}
A natural candidate invariant is the compensation of rainfall changes by irrigation, $\kappa$ (\cref{eq:kappa}). It is not universal (\cref{fig:comp}, \cref{tab:kappa}). In single-soil lattices the slope of $W$ on rain is $0.58\pm0.02$ in E11 ($\kappa=0.42$) and 0.33 in E12 and E14 ($\kappa=0.67$), whereas in E16 the pooled $\kappa$ across the ten soils is 0.159 (bootstrap 95\,\% interval 0.134--0.181) and the per-soil values range from $-0.014$ (sand) to 0.303 (coarse loam). Within E16, $\kappa$ falls from 0.193 to 0.123 as $E_p$ is raised from 0.3 to 1.2, from 0.188 to 0.129 with warming from $+2$ to $+6$\,\textdegree C, and from 0.211 to 0.107 with trigger suction from 50 to 2000\,cm, but rises from 0.079 to 0.241 with irrigation dose from 1.2 to 30\,mm. Compensation therefore depends on the scheduling rule and soil, and a single value of $\kappa$ should not be used in water-allocation arguments.

\begin{figure}[H]\centering
\includegraphics[width=\textwidth]{\figdir/p2/fig06_compensation.pdf}
\caption{Irrigation compensation $\kappa$ is not universal. (a) Per-soil $\kappa$ in E16 (paired $P=0.5\to1.5$) with the pooled estimate and bootstrap interval. (b) Slope of water input on rain in single-soil lattices (95\,\% intervals). (c) Mean $\kappa$ by factor level in E16.}
\label{fig:comp}
\end{figure}

\begin{table}[H]\centering\small
\caption{Compensation $\kappa=-\Delta I/\Delta P_r$ and the response of ET and drainage to rain, per soil in E16.}
\label{tab:kappa}
\begin{tabular}{lrrrr}
\toprule
Soil & $\kappa$ & $\Delta\mathit{ET}/\Delta P_r$ & $\Delta D/\Delta P_r$ & Pairs\\\midrule
S01 & 0.303 & 0.512 & 0.186 & 32\\
S02 & 0.164 & 0.346 & 0.495 & 8\\
S03 & $-0.014$ & 0.473 & 0.550 & 32\\
S04 & 0.084 & 0.473 & 0.453 & 32\\
S05 & 0.102 & 0.551 & 0.358 & 32\\
S06 & 0.108 & 0.619 & 0.281 & 32\\
S07 & 0.216 & 0.542 & 0.246 & 32\\
S08 & 0.090 & 0.613 & 0.308 & 32\\
S09 & 0.238 & 0.515 & 0.253 & 32\\
S10 & 0.302 & 0.477 & 0.237 & 32\\\bottomrule
\end{tabular}
\end{table}

%==========================================================

\subsection{The unattributed term and the recoverable share}\label{sec:auditlink}
The memory analysis measures an evaporation change by the non-transpiration term $E_{nt}=\mathit{ET}-T$, which carries the unattributed term $R$ with it. A user who read the reported soil evaporation $E_s$ alone would measure a different quantity. For the mulch experiment E7 the slope of drainage on $E_s$ was computed in place of the slope on $E_{nt}$, with the same regression that gives $\psi_D$ in the memory analysis (which reproduces the archived values for all ten soils; \cref{tab:auditlink}). The change of $E_s$ carried a median of 47\,\% of the change of $E_{nt}$ (range -0.40 to 0.56), the drainage share differed by more than 0.1 in 6 of 10 soils, and in 2 coarse soils (S03, S04) it changed sign (-0.93 and -0.35 against 0.71 and 0.69). The partition that the model reports therefore cannot be used to judge whether a saving is recoverable; the closed quantity $E_{nt}$ can. The analysis (\texttt{p1\_audit\_link.py}, output \texttt{results/p1\_audit\_link.json}) concerns one experiment of the model and is not a statement about measured soils.
\begin{table}[H]\centering\small
\caption{Drainage share of the change of evaporation in the mulch experiment E7 (20 combinations per soil), computed against the closed non-transpiration term $E_{nt}=\mathit{ET}-T$ (as in the memory analysis) and against the reported soil evaporation $E_s$ alone. $f_{E_s}$ is the share of the change of $E_{nt}$ that is carried by $E_s$.}\label{tab:auditlink}
\begin{tabular}{lrrr}\toprule
Soil & $\psi_D$ against $E_{nt}$ & $\psi_D$ against $E_s$ & $f_{E_s}$\\\midrule
S01 & 0.30 & 0.31 & 0.56\\
S02 & 0.61 & 0.38 & 0.46\\
S03 & 0.71 & -0.93 & -0.37\\
S04 & 0.69 & -0.35 & -0.40\\
S05 & 0.66 & 0.16 & 0.07\\
S06 & 0.52 & 0.41 & 0.48\\
S07 & 0.40 & 0.42 & 0.52\\
S08 & 0.55 & 0.36 & 0.39\\
S09 & 0.43 & 0.44 & 0.50\\
S10 & 0.41 & 0.45 & 0.50\\
\bottomrule\end{tabular}\end{table}

\section{Discussion of the memory analysis}
\label{sec:disc}
\subsection{Established findings and the operational definition of connectivity}
Within a mechanistic simulator, the fraction of an evaporative change that is recoverable depends on the soil in a way that is ordered by the pore-size-distribution parameter and by saturated conductivity, and the dependence is replicated under three different perturbations (antecedent moisture, evaporation factor, mulch). A simple rate-competition law explains the memory fate with a single parameter. Two climate drivers act in opposite directions on recoverability: warming increases the recoverable share (more of the extra evaporation is met by irrigation and transpiration rather than lost water) and a wetter climate reduces it. These are the findings the data support.

Pore connectivity is given here an operational definition in terms of water fluxes. In the Mualem conductivity function, connectivity enters through the exponent $\lambda$ and the tortuosity of the pore space \citep{Mualem1976,Ghanbarian2013}, and through $n$ and $K_s$; with $\lambda=0.5$ for all soils, the connectivity contrast among the ten soils is carried by the pore-size distribution and the saturated conductivity, and the memory fate measures its hydrological consequence directly. This is the useful reading: connectivity is expressed as a number, the share of released storage that leaves by drainage, that can be measured in a lysimeter without imaging the pore network. The natural complement is to pair the flux diagnostics with X-ray computed-tomography network metrics on the same soils, for example following \citet{Leuther2024}, so that the geometric and hydrological descriptions of connectivity can be related.

\subsection{Interpretation of why drainage competes with evaporation}
The soil's matrix acts as a reservoir with a conductance. When the profile is wetter than the field-capacity state, the unsaturated conductivity in coarse soils falls only slowly with suction, because a narrow pore-size distribution implies a rapid but late desaturation (large $n$), so the drainage flux stays high until the pores empty \citep{Hillel1998,Brady2016}. In fine soils, conductivity drops by orders of magnitude over the same suction range and the store is retained. When evaporation changes, the change in water remaining in the profile is shared between drainage (which responds on the timescale of $K$) and the irrigation or root uptake that responds on the timescale of the scheduler or crop. The parallel-sink form \cref{eq:rate} captures the competition between those timescales in its simplest form. The fitted $q^\ast$ being much larger than ET reflects that the relevant competition occurs during short, wet intervals after rain or irrigation, when conductivities are near saturation and drainage rates are of order of $K_s$ itself, and not over the annual mean.

The two signs are consistent with the same mechanism. Warming raises evaporative demand and dries the profile, so there is less water in the high-conductivity moisture range and less drainage to be displaced; extra evaporation then draws more from irrigation and transpiration. Rainfall has the opposite effect, keeping the profile wet and drainage high. That nitrogen shows no effect is a useful negative control because N influences yield and N leaching but is not coupled to the water partition in the model.

Where $\rho_R$ is low, in sands and loamy sands, a mulch saves much less irrigation water than its evaporation suppression suggests, because the saved water becomes deep drainage. Under those conditions more benefit may come from changing irrigation timing and depth than from reducing evaporation \citep{Fereres2007,Playan2006}. In contrast, in fine soils (S07, S09, S10; $\rho_R=0.64$--0.69) reducing evaporation translates into reduced irrigation need, supporting the use of residue mulches \citep{BalwinderSingh2011,Vuaille2024}. These conclusions are made for annual, 20-year means in a temperate climate and a one-dimensional model, and do not address field heterogeneity, preferential flow or the effects of mulch on infiltration and runoff \citep{Jarvis2007,Beven2013}. A companion software article (Paper~2) embeds a reduced form of these relationships in a decision-support tool and reports its agreement with the simulations.

\subsection{The competing rate as a physical scale}
The fitted competing rate $q^\ast=71.5$\,cm\,d$^{-1}$ (bootstrap 58--105, jackknife 66--80) is two orders of magnitude above mean evapotranspiration, and this is a finding in its own right. It identifies where the competition between drainage and withdrawal is decided: in the short wet intervals after rain or irrigation, when the profile is near saturation and drainage proceeds at rates of the order of $K_s$, not over the seasonal mean of uptake. The law $\phi_D=K_s/(K_s+q^\ast)$ then has two regimes with distinct consequences. For $K_s\ll q^\ast$ (the fine soils, down to $1.7$\,cm\,d$^{-1}$) the memory fate is proportional to $K_s$, so the uncertainty of a pedotransfer estimate of $K_s$ passes to $\phi_D$ in proportion; for $K_s\gtrsim q^\ast$ (the coarse soils, up to $713$\,cm\,d$^{-1}$) the fate saturates, so that further increases in conductivity change little. Both regimes are visible in the data, and the transition at $K_s=q^\ast$, where half of the released storage drains, is a prediction that can be located in a soil collection.

The law also compares favourably with the structural proxies that are available without conductivity data. In leave-one-soil-out validation on the eight soils with literature conductivities, the rate-competition law has an RMSE of 0.058, against 0.100 for $n$, 0.064 for the field-capacity saturation $S_{e,100}$ and 0.088 for the plant-available water, and 0.338 for the mean alone. On all ten soils the structural proxies give 0.097, 0.136 and 0.183 against 0.318 for the mean. The three proxies are collinear across the soil series, so their performance is similar and their combination adds little; $K_s$ is the axis that is not a function of $n$, and it is the one that the law uses. Because $K_s$ was taken from independent literature class means \citep{Carsel1988} and not from the simulator inputs, the law links an emergent property of the simulated soils to an externally measured property, so the agreement is not circular. Finally, $q^\ast$ is a quantity with the units of a flux, so the law yields a testable prediction for weighing lysimeters (Section~\ref{sec:disc}, testable predictions) that a structural index does not.

\subsection{Relation to earlier work}
Field and modelling studies of mulching document that evaporation falls, that the transpiration share of evapotranspiration rises, and that the response is stronger in finer-textured soils \citep{Zheng2021,Kang2024}; they report the size of the evaporation reduction and its partition, and not the share of the saved water that subsequently drains as a function of the soil's hydraulic properties, which is the quantity defined here. This analysis complements the soil-moisture memory literature \citep{Koster2001,Seneviratne2010} by quantifying not how long a moisture anomaly persists but where released storage goes. It extends the Budyko framework \citep{Budyko1974,Milly1994}, which treats annual partitioning as a function of climate aridity and (via a parameter) catchment properties, by showing that the sub-annual competition between drainage and withdrawal, which is not captured by the single Budyko parameter, is soil-controlled. It is consistent with the finding that hydrologic connectivity constrains the partitioning of terrestrial fluxes \citep{Good2015}, and with the efficiency paradox \citep{Grafton2018,Ward2008}, to which it adds a soil-scale mechanism: the saved water may be drainage. The need for structure-aware soil hydrology in land models \citep{Fatichi2020,Vereecken2016,Vereecken2019} is supported by the strong dependence of $\phi_D$ on the pore-size distribution.


Three testable predictions follow from the results. First, the drainage share of released antecedent storage should increase with saturated conductivity along the saturating form of the rate-competition law, so that a soil with $K_s$ equal to $q^{*}=71.5$\,cm\,d$^{-1}$ releases half of its stored water as drainage. Second, the recoverable share of an evaporation change should decrease in wetter climates and increase under warming, as found in nine of nine and ten of ten soils respectively. Third, the ranking of soils by $\phi_D$ should be preserved when the weather record is replaced by a different climate, because the ranking is set by hydraulic conductance and not by the weather. Each prediction can be tested with additional virtual experiments, and the first two with lysimeters in which drainage and evapotranspiration are measured on contrasting soils under imposed evaporation reduction by mulching.

In management terms, the results indicate that the saving obtained from mulching or from a lower evaporation factor is largely recoverable on fine-textured soils, where most of it appears as water that remains in the root zone, and is largely drainage on coarse soils, where it leaves the profile before the crop can use it. A mulch recommendation therefore requires a statement of the soil and of the climate, and the diagnostics defined here supply the quantities needed to make it. In land-surface and crop models the same diagnostics can be computed from a pair of runs, and they offer a compact check on whether a model's response to an evaporation change is plausible for the soil in question.


\paragraph{Relation of $\phi_D$ to the van~Genuchten parameter $n$.} The memory fate rank-orders with $n$ (Spearman $\rho=0.95$), as it should, since both derive from the pore-size distribution. They differ in kind. The parameter $n$ is a static, geometric descriptor, whereas $\phi_D$ is an operational quantity defined on water fluxes: the fraction of released antecedent storage that leaves as drainage, which depends on the withdrawal rate imposed by climate and crop as well as on the soil. The rate-competition law adds the saturated conductivity $K_s$, which is not a function of $n$, and maps soil properties onto a single coefficient that can be applied directly to an evaporation-saving claim. The contribution is therefore the definition, the mechanism and the transferability across experimental designs ($\rho=0.85$ and 0.89), with a prediction accuracy that is at least that of the structural proxies and, on the eight soils with conductivity, better than any of them.

\subsection{Drivers of soil hydraulic properties and the soil--plant--atmosphere continuum}
\label{sec:drivers}
The diagnostics defined here depend on soil hydraulic properties only through $K_s$ and the pore-size parameter $n$, but those properties are not fixed attributes of a soil. They are set by a small number of drivers and they respond to the plants and the atmosphere that the soil exchanges water with, so that the recoverability of an evaporation saving is a property of a coupled system (\cref{fig:drivers}, \cref{tab:drivers}).

\paragraph{Texture and pore-size distribution.} The primary control is particle-size distribution, which determines the pore-size distribution and therefore the retention curve and the conductivity function. Class-mean parameters and pedotransfer functions \citep{Carsel1988,Schaap2001,Saxton2006,Weber2024} encode this control, and it is the control carried by the ten soils of this ensemble (\cref{fig:drivers}a). Across them, $K_s$ spans more than two orders of magnitude (from 1.7 to 713\,cm\,d$^{-1}$ in the eight soils with literature values), and the rate-competition law \cref{eq:rate} maps this range onto a memory fate between 0.02 and 0.85.

\paragraph{Organic matter, aggregation and tillage.} Organic matter, biological activity and aggregation create structural porosity that raises $K_s$ and broadens the pore-size distribution relative to the texture-class value, and tillage and traffic change structure in space and in time \citep{Hillel1998,Brady2016,Strudley2008}. Because the saturated conductivity of structured soils can differ from the texture-class value by a factor of several, structural drivers move a soil along the curve of \cref{fig:drivers}a. Starting from $K_s=20$\,cm\,d$^{-1}$, for which \cref{eq:rate} gives $\phi_D=0.22$, a threefold increase in $K_s$ gives 0.46 and a threefold decrease gives 0.09. These values are implications of the fitted law and are not simulation results; the ensemble contains tillage and consolidation treatments (E5), but the memory fate was not computed for them in this paper.

\paragraph{Roots, macropores and non-equilibrium flow.} Root channels, earthworm burrows and shrinkage cracks conduct water preferentially and can decouple drainage from the matrix retention curve \citep{Jarvis2007,Beven1982}. A matrix-flow model, as used here, does not represent this pathway, so the drainage shares reported are those of the matrix alone; in a soil with active macropores the drainage share of released storage would be expected to be at least as large as the matrix value, and the recoverable share correspondingly smaller. This is a direction of change implied by the mechanism and not a measured effect. The recoverable share reported here should therefore be read as an upper bound for soils with active macropore flow.

\paragraph{Wetting--drying, freezing and surface condition.} Cycles of wetting and drying, and of freezing and thawing, rearrange pore structure and alter retention and conductivity \citep{Leuther2021}; residue and mulch cover change the surface hydraulic boundary and the infiltration and evaporation fluxes that act on it \citep{Leuther2026,Vuaille2024}. These drivers act on the same two properties, $K_s$ and $n$, and are therefore expressible in the diagnostics of this paper once their effect on the properties is known.

\paragraph{Coupling to plants and atmosphere.} Plants link the hydraulic properties to the atmosphere in both directions. Root distribution and uptake determine how much of the stored water is withdrawn before it can drain, which is the competing process in \cref{eq:rate}; evaporative demand, set by radiation, temperature and humidity, determines the withdrawal rate that drainage competes with \citep{Roderick2002,Brutsaert1998}; and the near-surface hydraulic conductivity at low suction controls how long the soil can sustain the evaporative demand at the atmospheric rate before the evaporation front recedes \citep{Or2013,Lehmann2008,Shokri2011}. The same soil hydraulic functions therefore regulate both sides of the competition: $K(h)$ near saturation governs drainage, and $K(h)$ at intermediate suction governs the supply of water to the surface and to roots. At larger scales, soil moisture feeds back on the surface energy balance and on precipitation in regions of strong coupling \citep{Entekhabi1996,Koster2004,Seneviratne2010}, and ecohydrological theory treats the soil moisture state as the integrator of these exchanges \citep{Rodriguez2004}. A recoverable share that depends on climate, as found here for warming and rainfall, is the soil-scale expression of this coupling: the climate changes the withdrawal rate that drainage competes with, and the resulting soil moisture state changes the evaporation and transpiration that follow.

\paragraph{Consequences for models.} Because structure and its dynamics modify hydraulic properties over seasons, representing them in land-surface and crop models is a recognised need \citep{Fatichi2020,Vereecken2016,Vereecken2019}. The relevant outcome here is practical: a recoverability diagnostic computed from $K_s$ and $n$ is only as stable as those properties, so the diagnostic should be evaluated for the structural state that management produces and not only for the texture-class value.

\paragraph{Change over time.} Hydraulic properties are not constants of a soil but state variables with their own time scales. Loosening by tillage raises the saturated conductivity and the macroporosity of the seedbed, and rainfall impact, wetting and traffic then re-consolidate and seal it over weeks to seasons \citep{Hillel1998,Brady2016,Strudley2008}; wetting--drying and freeze--thaw cycles rearrange pore structure on seasonal to annual scales \citep{Leuther2021,MirandaVelez2023}; and residue incorporation changes the hydraulic functions through the decomposition of the residue \citep{Leuther2026}. For a rate-competition law of the form of \cref{eq:rate}, any such change enters through $K_s(t)$, so the memory fate becomes a function of time, $\phi_D(t)=K_s(t)/[K_s(t)+q^\ast]$. \Cref{fig:timemgmt}a shows this for an illustrative case in which a threefold loosening of $K_s$ relaxes exponentially to the consolidated value with a rate $\tau_c$ of 0.005 or 0.02\,d$^{-1}$, two of the consolidation rates simulated in experiment E5 (0.001 to 0.02\,d$^{-1}$). The curves are implications of the fitted law under assumed multipliers and rates and are not simulation results. For the loam S06 ($K_s=25$\,cm\,d$^{-1}$, $\phi_D=0.26$ at the consolidated state), the memory fate is 0.51 immediately after loosening; with $\tau_c=0.02$\,d$^{-1}$ it has fallen to 0.32 after 90 days, whereas with $\tau_c=0.005$\,d$^{-1}$ it is still 0.44. The recoverable share of an evaporation saving therefore depends not only on the soil but on the time elapsed since its structure was last disturbed. The same mulch is therefore expected to lose a larger share of the water it saves to drainage in a freshly loosened, fast-draining seedbed than in the same soil late in the season; this direction of effect is a prediction to be tested, not a result.

\paragraph{Management practices.} Management acts on the same two properties through several routes. Tillage intensity and timing set the initial structural state and its decay; controlled traffic and reduced compaction keep $K_s$ closer to the structural value; organic amendments and residue retention increase aggregation and biological activity and, over years rather than seasons, shift the pore-size distribution \citep{Brady2016,Hillel1998}; surface mulch and residue cover modify the evaporation boundary, which is the quantity whose saving is being tracked here \citep{Vuaille2024,Leuther2026}; and irrigation method and trigger determine how often the profile is brought to the wet state in which drainage competes most strongly with uptake. A strength of the recoverable-share framing is that these practices can be compared on a common footing: each is expressed as a change in $K_s$ (or $n$) and read off the same curve. \Cref{fig:timemgmt}b shows the consequence of a threefold increase or decrease of $K_s$ for the eight soils with literature conductivities. The effect is strongly soil-dependent, because the law is concave on a logarithmic axis: in the loam S06, $\phi_D$ changes from 0.10 to 0.26 to 0.51 across the three cases, in the sandy S05 from 0.33 to 0.60 to 0.82, and in the silty clay loam S10 only from 0.008 to 0.023 to 0.066. A structural change of a given relative size matters most where the soil is intermediate, and little where the matrix is so tight that drainage is negligible whatever the management.

\paragraph{Soil-specific implications.} Three consequences follow for practice. First, the discount that should be applied to an evaporation-saving claim is a property of the soil--management--climate combination, not of the practice alone. Second, a soil that is close to the upper end of the curve (coarse, high $K_s$) gains little from structural improvement in this respect, whereas a soil in the steep middle of the curve is the most sensitive to compaction and to loosening. Third, because the loss fate depends on climate (rising with warming and falling with wetter conditions in the simulations), the same management has to be assessed under the climate of the site and under projected changes, not under the long-term mean only. None of these statements has been tested against field observations of structure over time in this paper.

\paragraph{Evidence from the simulation ensemble.} Two experiments of the Daisy ensemble bear directly on structure, time and mulching (\cref{fig:ens}); they are simulations and not field measurements. In run set E5, tillage and the rate $\tau_c$ at which the loosened soil consolidates were crossed for all ten soils, under rainfed and irrigated management. For the four finer soils with an effect (S06, S08, S09, S10) rainfed drainage rises with faster consolidation, by 1.1, 2.5, 2.1 and 1.5\,mm\,yr$^{-1}$ between $\tau_c=0.001$ and $0.02$\,d$^{-1}$, towards the no-till value (\cref{fig:ens}b,c); in S01--S05 drainage does not depend on $\tau_c$ at all, although tillage itself shifts it by $-1.8$ to $+5.3$\,mm\,yr$^{-1}$. The time dynamics therefore matter in the simulations only for the finer soils, and the effect is small against the texture contrast (rainfed drainage from 5 to 63\,mm\,yr$^{-1}$ across soils). Because the dynamic parameters of Daisy change retention and conductivity together, this response cannot be attributed to $K_s$ alone, and the multipliers of \cref{fig:timemgmt} remain assumptions; the ensemble indicates that, within the matrix-flow model, the structural term is second-order relative to texture. In run set E7 (strong against weak mulch, vapour-flux factor 0.3 against 0.9 at an interception capacity of 2\,mm), the evaporation saved ranges from 3 mm\,yr$^{-1}$ in the sand S03 to 54 mm\,yr$^{-1}$ in S01, and the extra drainage from 6 to 27\,mm\,yr$^{-1}$ (\cref{fig:ens}a); irrigation saved is near zero in S03--S05 and 43--49\,mm\,yr$^{-1}$ in S01, S07, S09 and S10, which is the same texture dependence as the shares $\psi_D^{E7}$ of \cref{tab:fates}. Measured $K_s(t)$ series and paired field plots remain the data needed to replace the assumptions.

\begin{figure}[H]\centering
\includegraphics[width=\textwidth]{\figdir/p1/fig14_ensemble_structure.pdf}
\caption{Simulation evidence on structure, time and mulching (Daisy ensemble, mean annual values). (a) Run set E7, strong against weak mulch (vapour-flux factor 0.3 against 0.9, capacity 2\,mm): evaporation saved, irrigation saved and extra drainage by soil. (b) Run set E5, rainfed drainage of four finer soils against the consolidation rate $\tau_c$ for conventional tillage (solid) and the no-till value (dotted). (c) Effect of tillage ($\tau_c=0.001$ against no-till) and of faster consolidation ($\tau_c=0.02$ against 0.001) on rainfed drainage, all ten soils.}
\label{fig:ens}
\end{figure}

\paragraph{Data that would test these predictions.} The predictions above are testable with three kinds of measurement: repeated infiltration or conductivity measurements across a season after tillage, which give $K_s(t)$ and its decay rate; paired drainage and evapotranspiration measurements under mulched and unmulched plots on contrasting soils, which give $\psi_D$ directly; and structure descriptors (bulk density, aggregate stability, macroporosity) that explain why $K_s$ departs from the texture-class value. Such data would allow the assumed multipliers and decay rates of \cref{fig:timemgmt} to be replaced by measured values, and would test whether the one-parameter law of \cref{eq:rate} holds for soils whose structure changes within the season.

\begin{figure}[H]\centering
\includegraphics[width=\textwidth]{\figdir/p1/fig13_time_management.pdf}
\caption{Implied change of the memory fate with soil structure over time and with management, from the fitted rate-competition law (\cref{eq:rate}, $q^\ast=71.5$\,cm\,d$^{-1}$). These are illustrative scenarios under assumed multipliers and rates, not simulation results. (a) A threefold loosening of $K_s$ relaxing exponentially to the consolidated value, for the sand S05, loam S06 and silty clay loam S10; solid lines $\tau_c=0.005$\,d$^{-1}$, dashed lines $\tau_c=0.02$\,d$^{-1}$, dotted lines the consolidated baselines. (b) Memory fate for the eight soils with literature $K_s$ at one third of, equal to and three times the baseline $K_s$.}
\label{fig:timemgmt}
\end{figure}

\begin{figure}[H]\centering
\includegraphics[width=\textwidth]{\figdir/p1/fig12_drivers.pdf}
\caption{Drivers of soil hydraulic properties and their coupling to the soil--plant--atmosphere continuum. (a) Memory fate $\phi_D$ as a function of saturated conductivity $K_s$: simulated values for the eight soils with literature $K_s$ (points) and the rate-competition law (line); arrows show the implied change for a threefold increase or decrease in $K_s$ from 20\,cm\,d$^{-1}$ (law implication, not simulation). (b) Conceptual feedback among drivers, hydraulic properties, water fates, soil moisture state, atmosphere and plants.}
\label{fig:drivers}
\end{figure}

\begin{table}[H]\centering\small
\caption{Drivers of soil hydraulic properties, their expected effect on $K_s$ and on the memory fate through \cref{eq:rate}, and their treatment in this study.}
\label{tab:drivers}
\begin{tabularx}{\textwidth}{>{\raggedright\arraybackslash}p{2.6cm}>{\raggedright\arraybackslash}X>{\raggedright\arraybackslash}p{2.3cm}>{\raggedright\arraybackslash}p{3.0cm}}
\toprule
Driver & Principal effect on hydraulic properties & Expected sign: $K_s$ / $\phi_D$ & Treatment here\\\midrule
Texture & sets pore-size distribution, retention and conductivity & coarser: $+$ / $+$ & ten soils; tested\\
Organic matter, aggregation & structural porosity, broader pore-size distribution & $+$ / $+$ & not varied\\
Tillage, consolidation & loosening raises, consolidation lowers conductivity & loosening: $+$ / $+$ & E5 in ensemble; $\phi_D$ not computed\\
Roots, macropores & preferential flow paths & $+$ / $+$ & not represented (matrix flow)\\
Wetting--drying, freeze--thaw & rearranges structure & sign depends on soil & not represented\\
Surface sealing, mulch & alters surface boundary, infiltration, evaporation & sealing: $-$ / $-$ & mulch tested (E7, E14)\\
Warming, rainfall & changes the withdrawal rate that drainage competes with & --- / drained share $\psi_D$: warming $-$, rain $+$ & E11 tested\\\bottomrule
\end{tabularx}
\end{table}

\subsection{Domain of validity and the experimental programme}
The conclusions are stated for the domain in which they were established, and each boundary of that domain defines an experiment.
\begin{enumerate}[nosep]
\item \emph{Soil series.} Ten soils span coarse to fine textures and more than two orders of magnitude in $K_s$. Soil-level statistics are reported with permutation tests and bootstrap intervals; the ordering of $\phi_D$ with $n$ and $K_s$ is firm ($\rho=0.95$ and 0.97), whereas the E16 association with $n$ (bootstrap $-0.06$ to 1.00) is indicative, and a larger soil library sampled from pedotransfer posteriors extends it.
\item \emph{Process representation.} The mechanisms are those of matrix flow in Daisy. Macropore bypass \citep{Jarvis2007,Beven1982} is absent, so the recoverable share is an upper bound for soils with active macropore flow, and replication in dual-permeability and Richards-based models such as HYDRUS, DSSAT, APSIM or AquaCrop \citep{Simunek2008,Simunek2016,Jones2003,Keating2003,Steduto2009} is the direct test of its structural generality.
\item \emph{Connectivity.} Defined through the Mualem function and the flux response (Section~\ref{sec:disc}.1).
\item \emph{Conductivity source.} $K_s$ is taken from independent literature class means for eight soils, which keeps the law non-circular; the remaining two soils are compared through $n$.
\item \emph{Initial-state design.} In E9, $h_0$ is the initial suction, so $\phi_D$ is the response to antecedent state, and its link to scheduling is established through the independent E16 and E7 experiments.
\item \emph{Time scale and climate.} Twenty-year means in one temperate climate; transfer to other climates is the second testable prediction below.
\item \emph{Completeness.} 10\,041 of 10\,190 runs (98.5\,\%) completed; the 149 failures include 13 in S03 within E9 and a large fraction of S02 within E16 (37 of 64 combinations valid), so the S02 loss fates are read as indicative.
\item \emph{Evidence status.} The conclusions are model-based and are formulated as quantitative predictions for lysimeter experiments; no field measurements enter the analysis.
\end{enumerate}

\Cref{tab:mgmt} translates the results into a qualitative guide. The recoverable fraction $\rho_R=1-\psi_D$ is quoted from E16 for the central climate; because \cref{fig:e11} shows that $\psi_D$ varies by an order of magnitude with rainfall and warming in a single soil, the table is a ranking aid and not a prediction.
\begin{table}[H]\centering\small
\caption{Qualitative reading of the diagnostics for management. $\rho_R$ from E16 (\cref{tab:fates}); ``mulch'' column from E7.}
\label{tab:mgmt}
\begin{tabularx}{\textwidth}{>{\raggedright\arraybackslash}p{3.3cm}>{\raggedright\arraybackslash}X>{\raggedright\arraybackslash}X>{\raggedright\arraybackslash}X}
\toprule
Soil group & Memory $\phi_D$ and $\rho_R$ & Evaporation-saving practice (e.g.\ mulch) & Where the saved water goes\\\midrule
Sand, loamy sand, sandy loam (S03--S05) & high $\phi_D$ (0.64--0.85); $\rho_R\le0.42$ & small irrigation benefit; $\psi_I^{E7}\approx0$ & mainly deep drainage\\
Loam, sandy clay loam, coarse sand (S06, S08, S02) & intermediate ($\approx0.26$--0.29) & moderate; $\psi_D^{E7}=0.52$--0.61 & split between drainage and irrigation\\
Silt loam, coarse loam, clay loam, silty clay loam (S07, S01, S09, S10) & low (0.02--0.15); $\rho_R\approx0.64$--0.71 & largest irrigation benefit; $\psi_I^{E7}=0.58$--0.63 & retained in root zone; irrigation and transpiration\\\bottomrule
\end{tabularx}
\end{table}

Land-surface schemes typically parameterise drainage with a single free-drainage lower boundary and a texture-based pedotransfer function \citep{Vereecken2019,Weber2024}. The strong dependence of $\phi_D$ on $K_s$ and $n$ suggests that, for evaporation-change questions, the uncertainty of the pedotransfer function enters the answer almost linearly through $K_s$ when $K_s\ll q^\ast$. Because pedotransfer functions have large uncertainty in $K_s$ (often an order of magnitude), the recoverable fraction of fine soils should be treated as uncertain by a comparable relative amount, while for sandy soils the saturating form of \cref{eq:rate} compresses the uncertainty. The absence of soil structure in most Earth system models \citep{Fatichi2020} means that the memory fate may be misrepresented precisely where structure is most important (fine soils, tillage).

(i) Extend the soil library by sampling parameters from pedotransfer-function posterior distributions and estimate the uncertainty of $q^\ast$. (ii) Repeat the E9, E16 and E7 designs in a dual-permeability model and in HYDRUS-1D \citep{Simunek2008}. (iii) Test the predictions of the Testable predictions subsection in lysimeters with controlled evaporation treatments and structure imaging. (iv) Examine climate beyond Taastrup, in particular Mediterranean climates where irrigation dominates the balance. (v) Link $\rho_R$ to economic and groundwater-recharge accounting \citep{Ward2008,Grafton2018}.

(i) In weighing lysimeters, the drainage response to a standard extra-evaporation treatment (shade or mulch) should scale as $K_s/(K_s+q^\ast)$ with $q^\ast$ of order 10$^1$--10$^2$\,cm\,d$^{-1}$. (ii) For two soils of equal $\TAW$ but different $n$, the one with the larger $n$ should show a smaller recoverable fraction. (iii) Under warming, the recoverable fraction should rise in all soils where irrigation is responsive. (iv) The recoverable fraction should be insensitive to nitrogen supply.

%==========================================================
\section{Discussion of the accounting audit}
\label{sec:discB}
\subsection{Contributions of a validated discovery protocol and a quantified accounting gap}
The paper contributes a discovery protocol in which every stage is tested against a known answer, and a first application that exposes an unclosed term in a widely used process model. The pipeline is blind to variable roles; it recovers the water balance as the dominant invariant, with a permutation null (smallest singular value 5.9 against 33.3 for permuted data), and this positive control is what licenses the interpretation of every other result. Applied to the evapotranspiration partition, the same pipeline quantifies a term that the reported components do not carry, $R\approx164$\,mm\,yr$^{-1}$ or about 30\,\% of ET, and a storage non-closure in two experiments (median 16\,\%). A compact three-term law then predicts ET across withheld experiments and withheld soils (LOEO $R^2=0.898$, LOSO 0.904) as accurately as gradient boosting (0.879 and 0.894). Quantifying a closure gap has a long record as a productive step in flux science: the energy-balance closure problem at eddy-covariance sites, in which turbulent fluxes typically fall 10--30\,\% short of available energy, was characterised across a network and then drove advances in measurement and in the interpretation of flux data \citep{Wilson2002,Foken2008}. The unattributed term $R$ has the same function for model-output accounting: it converts an implicit assumption (that reported components are the total) into a measured quantity that can be traced to a pathway. Stating the recovered invariant for what it is, the known mass balance returned by a blind method, is the credential that makes the new findings credible in a field where discovery claims are often overstated \citep{Rudin2019,Nearing2021}.

Evapotranspiration partitioning is a central topic, from global studies that infer transpiration fractions \citep{Jasechko2013,Schlesinger2014} to field studies of mulch effects \citep{BalwinderSingh2011}. If a model's reported ET components omit a third of the flux, then a partition fraction computed as $T/(T+E_s+E_c)$ is larger than $T/\mathit{ET}$ by a factor $1/(1-0.30)\approx1.4$ on average, an error comparable with the differences between management treatments that are routinely reported. The size of $R$ also helps to explain why the closure-derived and component-based estimates of soil evaporation differ. For users of Daisy and similar models, the practical recommendation is to compute ET as the model's own total and to report $R$ alongside the components, and to document, in the model output specification, exactly which processes enter the total \citep{Holbak2026}.

\subsection{A map of predictability across the flux budget}
The ET law works because ET is the large, smooth, sub-linear function of supply that the water-balance literature has described for decades, and because the model's ET is energy- and storage-limited \citep{Budyko1974,Milly1994,Zhang2001}. The drainage fraction is the small residual of two large numbers, it is strongly nonlinear in the wetness of the profile, and it depends on thresholds (field capacity, trigger) that a log-linear library cannot represent \citep{Rodriguez2004}; gradient boosting captures some of this and still only reaches $R^2\approx0.5$ in LOEO. This locates the information boundary: the content that drainage needs lies beyond the annual-flux library (e.g.\ intra-annual timing, event intensity, antecedent state). The memory analysis shows that antecedent state is indeed decisive for drainage in coarse soils.

The three-term law is competitive with gradient boosting under LOEO and LOSO, which supports the view that accuracy does not always require opacity \citep{Rudin2019}. The law does not extend to the regimes in which the physics it omits is first-order (initial-state experiments, dynamic hydraulic properties), and the gradient-boosting benchmark fails there too. This agreement is itself informative: it shows that the boundary is set by the information in the library and not by the model class, which tells a modeller which targets can be emulated by compact laws and which need timing and state variables.

Equation discovery in soil hydrology has so far targeted the governing partial differential equation, notably the Richards equation recovered from soil-moisture time series by group-sparse regression \citep{Song2022,Song2025}, and deep-learning studies have imposed water-balance closure as a constraint on forecasts \citep{Li2024}. The present paper is complementary and differs in its object and its emphasis: it works with annual flux budgets, tests whether a pipeline recovers a known balance before interpreting any other result, audits the closure of the partition that supplies its targets, and validates by withheld regimes. The present approach shares with SINDy-style discovery \citep{Brunton2016} the use of sparse regression over a library, but differs in two respects: it targets steady annual relations rather than dynamics, and it places the validation burden on leave-one-regime-out testing \citep{Roberts2017}. It shares with physics-informed learning \citep{Raissi2019,Karniadakis2021} the emphasis on conservation, but imposes it as a \emph{diagnostic} (closure residual) rather than as a constraint \citep{Beucler2021}; applying it as a constraint would have hidden $R$, which is the finding most relevant to model users. Related work on learning conserved quantities \citep{Liu2021,Iten2020} would be a natural extension to daily dynamics.

Emulators are used to make process-model knowledge available in fast decision tools \citep{Razavi2012,Asher2015}. The results here indicate which emulator targets are safe: ET and transpiration are predictable from compact laws; drainage and soil evaporation are not, and should be emulated with care or by closure.

A software-level audit would resolve the identity of $R$ directly. Four questions, in order of cost: (i) Which state variables are summed into the output called ``evapotranspiration'' in the log file used here, and which into the three reported components? (ii) Is evaporation from surface ponds, mulch and litter reported in a separate column that was not exported in the ensemble? (iii) Does the canopy-evaporation column represent interception only, excluding evaporation of water retained on residue? (iv) Is the storage term in the model's own water balance consistent with the exported depth and boundary definitions (for example, is the 1\,m drainage column the flux at the bottom of the root zone or the soil column)? If the answer to (ii) is yes, the unattributed term would be explained and the diagnostic here would have served its purpose as an audit.

Random $k$-fold cross-validation of the same models would have produced much higher scores, because runs inside an experiment are close neighbours in factor space. To illustrate the effect on the present data: a model that predicts $\ln\mathit{ET}$ from $\ln W$ alone has LOEO $R^2$ of 0.80 on the full data and 0.60 when E8 and E9 are removed. Leave-one-experiment-out removes entire designs, so the model must extrapolate to unseen combinations of factor ranges, which is exactly what a ``law'' is supposed to do \citep{Roberts2017}. The leave-one-soil-out scores answer a second question: would the law hold for a soil not used to fit it? That the two scores agree for ET (0.898 and 0.904) is evidence that the law is not driven by a single experiment or soil, but it is not evidence of validity in other climates. A nested evaluation, in which the forward selection was repeated inside each training set of 16 experiments and scored only on the held-out experiment, gives a pooled $R^2=0.895$, against 0.898 for the single-level score, and selects the same three terms in 15 of 17 outer folds (the two-term set $\ln W$, $n$ in the other two); the optimism from selecting on the evaluation score is therefore small for this law.

The sub-linear exponent $\partial\ln\mathit{ET}/\partial\ln W=0.646$ can be compared with the elasticity of Budyko-type curves. For a Budyko curve $\mathit{ET}/P=F(\Phi)$ with aridity index $\Phi=\mathit{PET}/P$, the elasticity of ET to $P$ lies between 0 (energy limit) and 1 (water limit) and is about 0.5--0.8 at intermediate aridity \citep{Milly1994,Zhang2001,Budyko1974}. The fitted value is within that range. Unlike the Budyko approach, which treats catchment properties via a single parameter, the present law includes a soil term ($n$): the fitted coefficient implies that a unit increase in $n$ reduces ET by 26\,\% ($1-e^{-0.299}$) at fixed $W$, which can be read as the dependence of the effective Budyko parameter on soil hydraulic uniformity. The law is not claimed to be a Budyko curve; the mapping would require PET as a variable, and PET is constant within the weather series used.


A sequential protocol follows from these results and is offered as a general template for discovery studies in hydrology. The first element is a positive control: the pipeline is applied to a system whose governing balance is known, and it must recover that balance with a null-referenced uncertainty, as the permutation test and bootstrap do here for the water balance \citep{Efron1993,Good2005}. The second is a closure audit of every partition used as a target, since a decomposition that omits a term will propagate that omission into every law fitted to it, as the unattributed term $R$ propagates into the transpiration and soil-evaporation fractions. The third is regime-wise validation, in which whole experiments, soils or climates are withheld, because random folds inside a designed ensemble measure interpolation between neighbouring factor combinations and not transfer to a new regime \citep{Roberts2017}. The fourth is a comparison with a flexible benchmark, reported with the cases in which the benchmark is better, so that the interpretability premium can be read directly \citep{Rudin2019,Friedman2001}. The fifth is a stated domain of validity for each law, expressed in the experiments and soils for which the leave-one-out score exceeds a stated threshold.

The unattributed term $R$ also defines a concrete verification programme for the process model. Three comparisons would identify it. First, an ablation in which individual surface and canopy evaporation pathways are switched off in turn would show which pathway removes $R$ from the balance. Second, a tracer-style accounting that records every term added to the evaporative flux in the model log, with its state variable of origin, would establish whether $R$ is a reporting convention or a physical pathway. Third, a comparison across structurally different models (Richards-based, capacity-based and radiation-driven) on the same designed ensemble would show whether a non-closing partition is a feature of one model family or a property of how evapotranspiration is partitioned in general \citep{Simunek2008,Jones2003,Keating2003,Steduto2009}. The experimental data that would constrain the physical interpretation are lysimeter and eddy-covariance records in which evaporation of intercepted and ponded water is measured separately from transpiration \citep{Good2015,Jasechko2013,Schlesinger2014}.

Beyond hydrology, the same sequence of controls is relevant to other settings in which a regression of a conserved quantity on its drivers is read as a discovered relation. Sparse regression over a library is attractive because the result is an equation that can be examined, dimensionally checked and compared with theory \citep{Brunton2016,Udrescu2020,Champion2019}. Its reliability, however, depends on the design of the data. In an ensemble built by crossing factors, several drivers are correlated by construction, and the selection path of a forward algorithm reflects that correlation; the leave-one-experiment-out score then serves two roles, as selection criterion and as validation, which is why the LOSO scores, which are less affected by selection, are reported alongside it. The practical consequence is that the coefficients of a discovered law should be reported with the validation regime in which they were obtained, and that a law should be described as discovered only for the regimes in which it is validated.

Finally, the results bear on the design of the emulators and decision tools that draw on process-model ensembles. A tool that reports evapotranspiration components should state whether they sum to the total; one that reports drainage should state whether it is simulated or derived by closure, and the evidence here shows that closure-derived drainage is accurate in sign and rank order but under-estimates large drainage. Emulators built on such ensembles should also be tested against held-out simulations \citep{Razavi2012,Asher2015,Wilkinson2016}.

\subsection{Domain of validity of the discovered laws}
Each law is reported with the regimes in which it was validated, and each boundary defines the next test.
\begin{enumerate}[nosep]
\item \emph{Process-model domain.} The results describe Daisy output. $R$ is a property of the model's accounting; its counterpart in nature is measured in lysimeters and eddy-covariance records that separate intercepted and ponded water from transpiration.
\item \emph{Failed runs.} The 149 excluded runs (87\,\% in S02) form a pattern that identifies them as solver failures (\cref{sec:failed}); re-running with tighter tolerances is the confirming step, and S02 is under-represented at large doses and extreme rain scales.
\item \emph{Model and weather.} One model and one climate; the laws are conditional on both, and the leave-one-regime-out protocol is designed to be repeated on further models and climates.
\item \emph{Selection.} Forward selection used the LOEO score; the nested evaluation, with selection repeated inside each training set, gives $R^2=0.895$ against 0.898 for the ET law, so the selection optimism is small. LOSO values are less affected by selection.
\item \emph{Dependence of runs.} Coefficient standard errors assume independent runs; cross-fold variation is the reported guide to uncertainty.
\item \emph{Soil covariates.} Ten soils; $n$ carries several correlated soil properties.
\item \emph{Identification of $R$.} Its surface-layer origin is inferred from its driver dependence, and the four-question audit of the model code in this section settles it.
\item \emph{Aggregation.} Twenty-year means average over seasonality and event structure; the memory analysis addresses antecedent state.
\item \emph{Benchmark reproducibility.} The gradient-boosting benchmark is run with a fixed random seed. Across unseeded repeats the LOSO score varied by up to 0.013, so differences below about 0.02 between a law and the benchmark are not interpreted.
\end{enumerate}

(1) Begin with a positive control with a known answer. (2) Use a permutation null for any decomposition-based invariant. (3) Validate across \emph{regimes} (experiments, soils, climates), not only across random folds. (4) Audit partitions for closure before using them as targets. (5) Compare with a flexible benchmark and report where the benchmark wins. (6) State a domain of validity for each law. (7) Publish failures.

%==========================================================
\section{What the two analyses say together}
The audit identifies which parts of the output of a process model can carry a conclusion, and the memory analysis shows what a conclusion built on them is worth. Three statements follow, each limited to the simulated system. First, an evaporation saving has to be judged against a closed quantity: $E_{nt}=\mathit{ET}-T$ closes by construction, whereas the reported components $E_s$ and $E_c$ leave out the unattributed term and give a different, in two soils opposite, answer (\cref{sec:auditlink}). Second, the soils in which drainage competes most with evaporation are the soils in which the partition is least settled, so the soil-specific discount of a saving cannot be replaced by a single factor. Third, the experiments that start away from equilibrium (E8, E9) are those in which storage does not close (16\,\%), and the memory fate was therefore estimated from the slope of drainage on storage release over many starting states and not from a single balance. None of this is evidence about measured soils: a lysimeter with controlled evaporation, measured drainage and a weighed mass balance would test the recoverable share and the unattributed term together.

\section{Conclusions}
Hydraulic memory, defined as the drainage share of released antecedent storage, ranges from 0.02 to 0.85 across ten simulated soils, follows the pore-size parameter ($\rho=0.95$) and saturated conductivity ($\rho=0.97$), and is captured out of sample by a one-parameter rate-competition law (LOSO RMSE 0.058). It transfers to independent evaporation-factor and mulch experiments ($\rho=0.85$, 0.89). The recoverable share of an evaporation change is soil-specific and climate-modified: lower in coarse soils (0.20 in loamy sand; none in sand), and shifted by $-0.16$ with warming and $+0.42$ with wetter climate in the drainage share. The accounting under these results was audited: a blind linear-invariant search recovers the water balance with every coefficient within 5\,\% of the exact value and flags an unattributed evapotranspiration term averaging 164\,mm\,yr$^{-1}$ (30\,\% of ET) and a storage non-closure of 16\,\% in the two initial-moisture experiments; a three-term law predicts evapotranspiration as well as gradient boosting in withheld experiments ($R^2=0.898$) and soils ($R^2=0.904$), and drainage resists both. Judged against the reported soil evaporation alone, the drainage share of an evaporation change would differ by more than 0.1 in six of ten soils and change sign in two. Together these give a transferable, mechanism-based recoverability diagnostic for ten soils and a protocol for deciding which outputs of a process model can carry a conclusion, with quantitative predictions that experiments controlling pore-network structure, antecedent state and evaporation can test.

\section*{Companion papers and release 1.1.0}
The simulation ensemble analysed here (17 experiments, 10\,190 runs, with the run time of every run, mean 148\,s, in the comprehensive summary of the repository) is also the basis of the decision layers of SWAFL Desktop 1.1.0. Paper~2 describes the engines built on it, Paper~3 the graphical Daisy workbench together with the bridge that connects Daisy to the decision layers, and Paper~4 the whole suite, including deterministic phenotyping, the field sketch, Sentinel-2 indices and the crop library. The results of both analyses of this paper were regenerated after the changes of release 1.1.0 and are unchanged.

\section*{Data and code availability}
\dataavail

\section*{Author contributions}
N.H.: conceptualisation, simulation, analysis, software, visualisation, writing.

\section*{Funding}
No funding was received for this work.

\section*{Competing interests}
The author declares no competing interests.

\section*{Acknowledgements}
The Daisy model is developed and maintained at the University of Copenhagen. The author thanks the maintainers of the open-source libraries cited.

\bibliography{refs}

\appendix
\section{Supplementary analyses}\label{app:hstar}
Replacing $K_s$ in \cref{eq:rate} by $K(h^\ast)$ of \cref{eq:mualem} with a fitted $h^\ast$ yields $h^\ast=6.7$\,cm and $q=15.4$\,cm\,d$^{-1}$ for the eight soils; in-sample RMSE is 0.027 and leave-one-soil-out RMSE 0.040 (S03: observed 0.846, in-sample 0.835, LOO 0.822). The variant has one more fitted parameter, uses the same eight soils, and is mildly better than the $K_s$ form; the choice between them cannot be resolved with eight points. The simpler form is therefore presented in the main text.

The 64 combinations per soil are written \texttt{C\#\_P\#\_T\#\_h\#\_D\#\_N\#\_Ep\#}. Valid combinations: 64 in nine soils and 37 in S02, so $\psi_D$ for S02 is estimated from fewer pairs; paired differences use only complete pairs.

\subsection{Mass-balance closure}
Median $|\Delta S|/W$ is 0.9\,\% over all valid runs; excluding E8 and E9 it is 0.65\,\% (95th percentile 2.3\,\%); for E8/E9 it is 16\,\%, as expected for experiments that deliberately start away from equilibrium. Fate shares in \cref{eq:phi} sum to 1 to within rounding.

\cref{fig:rec} shows the recoverable share $1-\psi_D$ for E16 and E7. Because E16 and E7 differ in how evaporation is changed, the two estimates differ in the order of $\pm0.2$ except in the sandy soils.
\begin{figure}[H]\centering
\includegraphics[width=0.8\textwidth]{\figdir/p1/fig09_recoverable.pdf}
\caption{Recoverable fraction $1-\psi_D$ by soil for the E16 and E7 contrasts.}
\label{fig:rec}
\end{figure}


\begin{table}[H]\centering\scriptsize
\caption{E9 reference values at $P=1$. storage release $\bar R$ (mm\,yr$^{-1}$), drainage $D$ and transpiration $T$ for wet ($h_0=30$ or 100\,cm) and dry ($h_0=1500$ or 3000\,cm) starts. Negative $D$ denotes net capillary rise across the 1\,m boundary; ``--'' marks a failed run (S03 at 3000\,cm).}
\label{tab:e9}
\begin{tabular}{llrrrrrrrr}
\toprule
Soil & Texture & $\bar R_{30}$ & $\bar R_{100}$ & $\bar R_{1500}$ & $\bar R_{3000}$ & $D_{100}$ & $D_{3000}$ & $T_{100}$ & $T_{3000}$\\\midrule
S01 & Coarse loam & 236.2 & 69.7 & -13.1 & -23.1 & -2 & -1 & 80 & 24\\
S02 & Coarse sand & 214.8 & 106.9 & 2.8 & -6.1 & 23 & 26 & 88 & 29\\
S03 & Sand & 3.2 & -24.1 & -28.2 & -- & 66 & -- & 26 & --\\
S04 & Loamy sand & 21.3 & -27.5 & -40.7 & -41.0 & 47 & 40 & 25 & 20\\
S05 & Sandy loam & 72.6 & -20.8 & -62.1 & -63.3 & 18 & 2 & 31 & 15\\
S06 & Loam & 197.3 & 84.6 & -47.2 & -51.8 & 9 & -0 & 116 & 13\\
S07 & Silt loam & 258.2 & 177.2 & -31.0 & -41.8 & 10 & 0 & 190 & 16\\
S08 & Sandy clay loam & 117.6 & 5.3 & -51.7 & -55.4 & 7 & -0 & 50 & 15\\
S09 & Clay loam & 192.9 & 139.9 & -27.5 & -38.7 & 4 & 0 & 167 & 21\\
S10 & Silty clay loam & 189.5 & 158.2 & -14.8 & -29.6 & 0 & 0 & 186 & 26\\
\bottomrule\end{tabular}\end{table}

\begin{table}[H]\centering\small
\caption{Closure of the loss fates. $\Delta\Ent$: mean paired change in non-transpirational evapotranspiration for $E_p=0.3\to1.2$ (mm\,yr$^{-1}$); $\psi_D,\psi_I,\psi_T$: shares; $\Sigma$: their sum (the remainder is storage and runoff); span E7: range of $\Ent$ across the mulch lattice (mm\,yr$^{-1}$).}
\label{tab:closure}
\begin{tabular}{lrrrrrrrr}
\toprule
Soil & $\Delta\Ent$ & $\psi_D$ & $\psi_I$ & $\psi_T$ & $\Sigma$ & span E7 & $\psi_D^{E7}$ & $\psi_I^{E7}$\\\midrule
S01 & 218 & 0.295 & 0.619 & 0.111 & 1.025 & 167 & 0.298 & 0.583\\
S02 & 144 & 0.193 & 0.658 & 0.157 & 1.008 & 86 & 0.607 & 0.339\\
S03 & 100 & 1.210 & -0.384 & 0.178 & 1.004 & 59 & 0.706 & 0.000\\
S04 & 115 & 0.799 & 0.219 & -0.007 & 1.011 & 58 & 0.687 & -0.023\\
S05 & 141 & 0.582 & 0.415 & 0.015 & 1.012 & 74 & 0.663 & -0.027\\
S06 & 191 & 0.469 & 0.433 & 0.107 & 1.009 & 122 & 0.521 & 0.379\\
S07 & 218 & 0.362 & 0.528 & 0.117 & 1.007 & 156 & 0.396 & 0.605\\
S08 & 152 & 0.533 & 0.386 & 0.091 & 1.010 & 111 & 0.549 & 0.332\\
S09 & 217 & 0.312 & 0.645 & 0.048 & 1.005 & 152 & 0.426 & 0.584\\
S10 & 222 & 0.310 & 0.838 & -0.017 & 1.131 & 160 & 0.414 & 0.629\\
\bottomrule
\end{tabular}
\end{table}

\section{Reproducibility map of the memory analysis}\label{app:repro}
\begin{table}[H]\centering\footnotesize
\caption{Result-to-script map (scripts in \texttt{Manuscripts/reproducibility/analysis/}, outputs in the work directory \texttt{out/}).}
\begin{tabularx}{\textwidth}{XX}
\toprule
Result & Script / output\\\midrule
Ensemble table and run counts & \texttt{make\_data.py} $\to$ \texttt{data/daisy\_ensemble.csv.gz}\\
Memory and loss fates (Tables \ref{tab:fates}) & \texttt{p1\_core.py} $\to$ \texttt{p1\_core.json}, \texttt{p1\_table.csv}\\
Rank statistics (Table \ref{tab:rank}) & \texttt{p1\_stats.py} $\to$ \texttt{p1\_table2.csv}; \texttt{p1\_figs\_d.py} $\to$ \texttt{p1\_rank.json}\\
Rate-competition law: $q^\ast$, $R^2$, bootstrap and leave-one-out intervals, LOSO errors (Fig.\ \ref{fig:mech}, Table \ref{tab:loso}) & \texttt{p1\_rate\_law.py} $\to$ \texttt{p1\_rate\_law.json}\\
Earlier two-parameter variant; logistic-in-$n$ LOO RMSE & \texttt{p1\_mech.py} $\to$ \texttt{p1\_mech.json}\\
Modifiers (Table \ref{tab:mod}) & \texttt{p1\_figs\_c.py} $\to$ \texttt{p1\_modifiers.json}\\
Figures & \texttt{p1\_figs\_a--e.py}, \texttt{p1\_fig12--14*.py}, \texttt{p1\_fig1\_new.py}, \texttt{figstyle.py}\\
All steps, checks and limits & \texttt{run\_all.sh}, \texttt{compare\_results.py}, \texttt{audit\_numbers.py}, \texttt{README\_REPRODUCE.md}\\\bottomrule
\end{tabularx}
\end{table}

\section{Supplementary information}
The analysis was exploratory and was not pre-registered. The main analytical decisions were: the four controls (permutation null, bootstrap, LOSO/LOEO, gradient-boosting benchmark); the forward-selection stopping threshold (0.003); log-targets; the exclusion of E8 and E9 as a sensitivity analysis for the invariant search and for closure-derived drainage; and the treatment of $R$ as a remainder. The candidate library was assembled by the author with knowledge of the ensemble and of the hydraulic-memory results (for example, the choice to include $n$, $\ln\alpha$ and $\ln W/\TAW$), which increases researcher degrees of freedom; readers should treat LOEO and LOSO scores as indicative rather than as confirmatory error rates.

The seven-flux search also returns weights on transpiration, soil evaporation and canopy evaporation in its smallest direction (rain 0.896, irrigation 0.834, ET $-1$, transpiration 0.367, soil evaporation 0.121, canopy evaporation $-0.122$, drainage $-0.933$), which is not a clean invariant: because $T+E_s+E_c+R=\mathit{ET}$ and $R$ is not a column, the seven-column matrix has an additional near-dependence that mixes components. This is why the four-flux system is the clean positive control.

\begin{table}[H]\centering\small
\caption{Columns of the ensemble used in this paper.}
\begin{tabularx}{\textwidth}{llX}
\toprule
Column & Unit & Definition\\\midrule
\texttt{experiment}, \texttt{treatment}, \texttt{run\_id} & -- & experiment code, treatment string encoding factor levels, unique run identifier\\
\texttt{soil} & -- & S01--S10\\
\texttt{rain\_mm}, \texttt{irrigation\_mm} & mm\,yr$^{-1}$ & 20-year mean annual rainfall and irrigation\\
\texttt{ET\_mm} & mm\,yr$^{-1}$ & model evapotranspiration (total)\\
\texttt{transpiration\_mm}, \texttt{soil\_evap\_mm}, \texttt{canopy\_evap\_mm} & mm\,yr$^{-1}$ & reported components of ET\\
\texttt{drainage\_1m\_mm} & mm\,yr$^{-1}$ & flux across the 1\,m boundary (negative = net rise)\\
\texttt{water\_stress\_d} & d\,yr$^{-1}$ & days of water stress\\
\texttt{yield\_Mg\_ha} & Mg\,ha$^{-1}$ & grain/biomass yield\\
\texttt{N\_fertiliser\_kg\_ha}, \texttt{N\_leached\_kg\_ha} & kg\,ha$^{-1}$ & nitrogen applied and leached\\
\texttt{IE\_kg\_m3}, \texttt{DPF}, \texttt{NLF} & -- & irrigation efficiency, drainage fraction $D/W$, N-leaching fraction (definitional; excluded as targets)\\
\texttt{P}, \texttt{dT}, \texttt{Ep}, \texttt{trig}, \texttt{z}, \texttt{amt\_mm}, \texttt{vff}, \texttt{cap}, \texttt{Nrate} & various & factor levels (rain scale, warming, evaporation factor, trigger, sensor depth, dose, bare fraction, mulch capacity, N rate)\\
\bottomrule
\end{tabularx}
\end{table}

DPF$=D/W$ is exact by construction and its ``prediction'' is therefore the same task as predicting $D$; NLF and IE are ratios of other columns. They are not used as independent evidence in any result.

\section{Reproducibility map of the accounting audit}\label{app:reproB}
\begin{table}[H]\centering\footnotesize
\caption{Result-to-script map.}
\begin{tabularx}{\textwidth}{XX}
\toprule
Result & Script / output\\\midrule
Ensemble table and run counts & \texttt{make\_data.py} $\to$ \texttt{data/daisy\_ensemble.csv.gz}\\
SVD, permutation null, closure (\cref{eq:law}, \cref{fig:invariants}) & \texttt{p2\_a.py}, \texttt{p2\_laws.py} $\to$ \texttt{p2\_a.json}, \texttt{p2\_laws.json}\\
Unattributed term $R$ (\cref{eq:Rlaw}, \cref{fig:hidden}) & \texttt{p2\_R.py} $\to$ \texttt{p2\_R.json}\\
clr-PCA (\cref{fig:manifold}, \cref{tab:clipsens}) & \texttt{p2\_a.py}, \texttt{p2\_figs.py}, \texttt{p2\_fig4.py} $\to$ \texttt{p2\_pca.json}, \texttt{p2\_pca\_sens.json}\\
Sparse laws (\cref{tab:laws}, \cref{eq:ETlaw}) & \texttt{p2\_et.py}, \texttt{p2\_dpf.py}, \texttt{p2\_figs2.py}, \texttt{p2\_nested.py} $\to$ \texttt{p2\_et.json}, \texttt{p2\_dpf.json}, \texttt{p2\_etlaw.json}, \texttt{p2\_nested.json}; exclusion sensitivity \texttt{p2\_s02\_excl.py} $\to$ \texttt{p2\_s02\_excl.json}\\
Cluster bootstrap, law without E8 and E9, descriptive statistics (Table~\ref{tab:desc}) and $R$ by experiment (\cref{tab:R}) & \texttt{p2\_robust.py} $\to$ \texttt{p2\_robust.json}\\
Compensation (\cref{fig:comp}) & \texttt{p2\_laws2.py} $\to$ \texttt{p2\_laws2.json}\\
Figures & \texttt{p2\_figs.py}, \texttt{p2\_figs2.py}, \texttt{p2\_fig4.py}, \texttt{p2\_fig6.py}, \texttt{p2\_fig7.py}, \texttt{p2\_figfail.py}, \texttt{p2\_fig1\_new.py}\\
All steps, checks and limits & \texttt{run\_all.sh}, \texttt{compare\_results.py}, \texttt{audit\_numbers.py}, \texttt{README\_REPRODUCE.md}\\\bottomrule
\end{tabularx}
\end{table}

\end{document}
